% https://habr.com/ru/companies/ruvds/articles/574352/
% https://guides.nyu.edu/LaTeX/sample-document
% https://ru.stackoverflow.com/questions/222769/latex-как-быть-с-русским-текстом/757767#757767

\documentclass{article}
% без этой строчки (модуль cmap) не будет работать поиск внутри документа и копирование русского текста:
\usepackage{cmap}
\usepackage[utf8]{inputenc}
\usepackage[english,russian]{babel}

% разбить длинные адреса URL, которые не умещаются в строку
% https://tex.stackexchange.com/questions/54946/how-to-break-a-long-url
\usepackage[hyphens]{url}

% использовать такой же шрифт в адресах \url, как в основном тексте
% https://tex.stackexchange.com/questions/261434/changing-url-font
\urlstyle{same}

% слова с дефисом (здесь в основном - "аппаратно-программная") по умолчанию не переносятся автоматом и вылезают за пределы строки
% чтобы этого не происходило, следует заменить дефис на команду "=
% https://www.linux.org.ru/forum/general/11119398
% http://www.inp.nsk.su/~baldin/LaTeX/ctex.pdf
% https://stackoverflow.com/questions/2193307/how-do-i-get-latex-to-hyphenate-a-word-that-contains-a-dash
% в текущем документе работает без подключения дополнительных модулей (в смысле, необходимые модули уже и так подключены выше)
% заменяю только в тех случаях, где после экспорта в pdf строки, действительно, вылезают за границы
% (в консоли вывода pdflatex в этом случае будет предупреждение "Overfull \hbox")

% Если после двойных кавычек идет пробел, при генерации pdf пробел уничтожается:
% например, в исходнике 'myth," he', превратится в pdf в 'myth,"he'
% В тексте использую кавычки ёлочки «»,
% если кавычки внутри кавычек (в цитатах), двойные обычные кавычки заменяю на лапки „“ или другие верхние двойные кавычки “”
% (лапки с нижней кавычкой нужно использовать в русских текстах, а две верхние кавычки - в английских)
% https://habr.com/ru/articles/865866/

% \raggedright внутри \section - отключить перенос слов внутри заголовка
% https://www.linux.org.ru/forum/general/11676660

\usepackage{graphicx}

% нужно включить, чтобы движок понимал символы евро € и умножение ×
% (символ фунта £ он понимает и так)
% но чтобы символ отображался, нужен еще подходящий шрифт
\usepackage{textcomp}

% шрифт по умолчанию не поддерживает символ евро (точнее, он рисует что-то похожее, но не то, что применяется сейчас)
% шрифт с поддержкой символа евро €, похожий на шрифт по умолчанию
% проблема с lmodern в том, что с ним перестает работать возможность выделения слов курсивом с \textit
% https://tex.stackexchange.com/questions/336155/italicization-with-textit-does-not-work#336158
% https://tex.stackexchange.com/questions/82116/how-to-display-italics-with-lmodern-package
\usepackage{lmodern}

% другие шрифты с поддержкой символа евро € и курсива одновременно
% можно использовать их, но мне не нравится их начертание, хочется оставить вариант максимально похожий на дефолт
%\usepackage{dejavu}
%\usepackage[scaled=0.9]{DejaVuSerif}
%\usepackage[scaled=0.8]{DejaVuSerif}
%\usepackage{libertine}
%\usepackage{paratype}

% использовать шрифт с поддержкой курсива для блоков \textit (блок \textit переопределяем)
% https://latexref.xyz/_005cnewcommand-_0026-_005crenewcommand.html
% https://tex.stackexchange.com/questions/47351/can-i-redefine-a-command-to-contain-itself
% это нужно для шрифта lmodern, если использовать шрифты с поддержкой символа евро € и курсива, то этот хак можно выключить
% 'cmr' - судя по всему, имя дефолтного шрифта, у него есть курсив, но нет символа евро €
% у шрифта lmodern есть символ €, но нет курсива
% для основного текста будем использовать lmodern, для курсива - 'cmr'
% (с таким вариантом символ € внутри блока \textit размещать нельзя - он его нарисует по-другому)
% внутренее имя шрифта 'cmr' узнал из ворнингов pdflatex при генерации документа с шрифтом без курсива
% кой-какие подробности про выбор шрифта для блока текста и внутренние имена шрифтов:
% % https://tex.stackexchange.com/questions/25249/how-do-i-use-a-particular-font-for-a-small-section-of-text-in-my-document
\let\textitorig\textit
\renewcommand{\textit}[1]{{\fontfamily{cmr}\selectfont \textitorig{#1}}}

% аналогично - для жирного текста в блоке \textbf
\let\textbforig\textbf
\renewcommand{\textbf}[1]{{\fontfamily{cmr}\selectfont \textbforig{#1}}}

% уменьшим отступ для \quad (по умолчанию \quad - это \hspace{1em})
\renewcommand{\quad}{\hspace{0.5em}}

% для блока align для многострочных формул
\usepackage{amsmath}

\title{Оценка эффективности цифровой трансформации предприятий и организаций}
\author{А.~Е.~Моисеев}
\date{cентябрь 2025}
%\date \today

\begin{document}

\maketitle

\begin{center}
\textit{на правах рукописи}
\end{center}

Методы оценки эффективности цифровизации компаний могут включать как субъективные экспертные, так и объективные метрики оценки уровня проникновения цифровых технологий в текущую деятельность компании и влияния, которые они на нее оказывают.

К примеру, показатели уровня цифровой зрелости, представленные в результатах исследования SAP и Deloitte «Цифровая зрелость российских компаний» в 2021, вычислены на основе опросов руководителей компаний \cite{sapDeloitteSurvey_eg_2021, sapDeloitteSurvey_deloitte_2021}. На основе результатов ответов на вопросы интервью в исследовании сформирован набор балльных метрик со значениями на непрерывной шкале от 1 до 5. Среди метрик, которые, по мнению авторов, оказывают прямое влияние на бизнес-показатели компании, перечислены «Клиентский опыт», «Клиентские инсайты», «Управление изменениями», «Интегрированное управление сервисом», «Онлайн-аналитика».

Индекс цифрового ускорения (Digital Acceleration Index — DAI), рассчитанный специалистами аналитической компании BCG в 2021, вычисляется также на основе интервью с руководителями, которые сами определяют степень «цифровой зрелости» своих компаний \cite{bcgSurvey2021}. В исследовании среди результатов цифровизации отмечают в том числе повышение производительности труда, к примеру, увеличение добычи или повышение эффективности применения производственного оборудования.

Индекс готовности к цифровому обществу (Network Readiness Index — NRI), рассчитываемый ежегодно Институтом Портулан (Portulans Institute) \cite{networkReadinessIndex2024}, формируется из нескольких численных метрик, вычисляемых для отдельных стран. Авторы методики определяют степень влияния цифровизации на экономику страны через значения четырех параметров: количество поданных патентов в области информационных и коммуникационных технологий (information and communications technology — ICT) на миллион человек, размер внутреннего рынка, доля экспорта высокотехнологичных информационных и коммуникационных продуктов в общем экспорте, уровень преобладания в стране элементов т.~н. гиг-экономики (экономика, в которой доминирует частичная занятость и фриланс, в современных российских реалиях — т.~н. платформенная занятость, самозанятость, трудоустройство по контрактам гражданско-правового характера ГПХ вместо устройства по трудовому кодексу ТК). Метрика готовности к гиг-экономике для каждой страны вычисляется ежегодно на основе результатов опроса участников Всемирного экономического форума.

В ноябре 2024 ВШЭ выпустила очередную редакцию ежегодного статистического сборника «Индикаторы цифровой экономики» \cite{hseDigitalStats2024}, в котором привела около ста восьмидесяти всевозможных показателей. Предложенная в сборнике модель статистического измерения цифровой экономики включает модуль «эффекты цифровизации», определяемый через «эффекты внедрения и использования цифровых технологий» и «экспорт цифровых технологий, товаров и услуг, связанных с цифровыми технологиями». Среди приведенных показателей широко представлены и детализированы затраты государства и организаций на цифровые технологии и исследования в области телекоммуникаций, статистика по имеющейся в наличии и применяемой инфраструктуре, применяемым технологиям, подготовке кадров и грамотности населения в сфере информационно-коммуникационных технологий (ИКТ) и т.~п. Детализированы деятельность организаций сектора ИКТ, сектор ИКТ как таковой, статистика импорта и экспорта продуктов ИКТ, деятельность сегмента контента и СМИ. Таким образом, в сборнике широко представлены структура применяемых технологий, структура и размер затрат на них. Эффекты влияния технологий на экономику за пределами сектора ИКТ, кроме самого факта их применения, представлены в меньшей степени. Среди имеющихся в наличии показателей «эффекты внедрения и использования цифровых технологий» отражены, к примеру, данными по оценке населением и организациями качества предоставленных электронных государственных и муниципальных услуг.

Росстат с 2002 года (на тот момент — Госкомстат) собирает данные о проникновении цифровых технологий в организации по форме 3-ИНФОРМ \cite{rosstat_prikaz332_31072024_3Inform}, в форму регулярно от года к году вносятся изменения и дополнения. Форма включает разделы с перечнем технологий, которые применяет организация, целями их применения, детализацией расходов организации на покупку и внедрение применяемых цифровых технологий и т.~п. Раздел с оценкой результатов внедрения в 2025 году включает набор показателей таких, как «Рост доходов», «Снижение себестоимости продукции», «Повышение качества продукции», «Повышение качества / эффективности производственных процессов», «Повышение качества / эффективности бизнес-процессов», «Рост производительности труда», «Сокращение сроков вывода продукции / услуг на рынок», «Рост числа клиентов», «Снижение численности работников», «Повышение безопасности труда». Предусмотренные варианты ответа на каждую из позиций — «да», «нет», «затрудняюсь ответить», — без детализации в конкретных числах. Раздел с информацией о применении электронных продаж существует с первой редакции формы 2002 года \cite{rosstat_postanov158_25072002_3Inform, rosstat_postanov39_04032003_3Inform}, позволяет оценить долю продаж, производимых организацией через Интернет, в процентах от общих продаж по другим каналам.

Показатели, отражающие расходы организаций такие, как собираемые по форме 3-ИНФОРМ, позволяют оценить рынок цифровых технологий — детализировать перечень применяемых технологий, а также его объем в денежном выражении. Организации несут издержки на покупку, разработку и внедрение цифровых технологий для того, чтобы достигать с их помощью экономических эффектов. Вполне очевидно, что при хорошей стратегии внедрения положительный эффект от внедрения превысит первоначально понесенные издержки. Таким образом, известные статистические данные с объемом расходов на цифровые технологии можно использовать для косвенной оценки размера экономического эффекта цифровизации — он должен быть, как минимум, не меньше понесенных издержек. В действительности, нет никакой гарантии, что это, действительно, так. Для более точной оценки эффективности внедрения цифровых технологий хорошо иметь в распоряжении данные, которые отражают результаты их внедрения в производственный процесс.

\section*{\raggedright{Вычисление экономических эффектов внедрения технологии на основе оценки производимой предприятием добавленной стоимости}}

Прямой непосредственный эффект внедрения новой технологии в производственный процесс — повышение производительности труда в точке внедрения. Организация экономического сотрудничества и развития (ОЭСР, OECD) и Международная организация труда (МОТ, ILO) определяют производительность труда в денежном выражении в масштабах экономики государства как отношение Валового внутреннего продукта (ВВП) к количеству примененного при его создании труда — количеству занятых на производстве продукта рабочих или количеству затраченных при производстве продукта человеко-часов \cite{oecdLabourProductivityManual2001, iloLabourProductivityKILM2016}.

В России Росстат рассчитывает «Индекс производительности труда», опираясь на количество фактически отработанных человеко-часов (приказ Росстата от 28 апреля 2018 года N~274 \cite{rosstat_prikaz274_28042018_LabourProductivityIndex}), ссылается на методику ОЭСР. Итоговое значение индекса представляет собой изменение метрики производительности труда в текущем году по сравнению с предыдущим годом в процентах \cite{rosstatLabourProductivityIndexData2025}. Минэкономразвития в рамках национального проекта «Эффективная и конкурентная экономика» рассчитывает показатель «Рост производительности труда», опираясь на значения показателя «Индекс производительности труда», рассчитываемого Росстатом (приказ Министерства экономического развития Российской Федерации от 23 декабря 2024 г. N~817 \cite{mineco_prikaz817_23122024_LabourProductivityGrowth}).

В то же время Росстат рассчитывает показатель «Динамика производительности труда», беря за основу производительность в пересчете на численность занятых на производстве рабочих (приказ Росстата от 16 февраля 2018 г. N~76 \cite{rosstat_prikaz76_16022018_LabourProductivityDynamics}). В рамках национального проекта «Производительность труда и поддержка занятости» Минэкономразвития использует аналогичный подход к расчету производительности, определяя её для отдельного предприятия в денежном выражении как отношение объема добавленной стоимости, произведенной на предприятии, к количеству сотрудников, задействованных при ее производстве (приказ Минэкономразвития от 28 декабря 2018 года N~748 \cite{mineco_prikaz748_28022018_LabourProductivityMethod}).

Имея в распоряжении данные со значениями индекса производительности и изменением его уровня такие, какие считает Минэкомразвития и публикует Росстат, удобно отслеживать динамику изменения производительности в экономике, отрасли или на отдельном предприятии, сравнивать производительность разных предприятий или одного и того же предприятии в разные моменты времени. Для того, чтобы сравнить результат повышения производительности с предварительными расходами на мероприятия, благодаря которым оно достигнуто, хорошо иметь в распоряжении абсолютное значение предварительных расходов на проведение этих мероприятий, а также оценить абсолютное значение реализованного с их помощью экономического эффекта.

Пусть $V_1$ — добавленная стоимость, произведенная на предприятии за некоторый период до внедрения новой технологии, при этом на производстве для её создания было задействовано $n_1$ рабочих. После внедрения новой технологии за аналогичный период времени была произведена стоимость $V_2$ трудом $n_2$ рабочих.

В том случае, если количество рабочих на производстве не изменилось ($n_2 = n_1$), прямой экономический эффект $E$ составит приращение добавленной стоимости \cite{vedomostiRussianIndustryEfficiency2024, minecoLabourProductivityProjectPresentation2023} — добавочную добавленную стоимость, — его можно оценить как:

$$E = (V_2 - V_1)$$

В том случае, если эффект внедрения $E$ больше нуля, производительность труда повысилась, в том случае, если $E$ равен нулю, производительность осталась без изменений, если $E$ меньше нуля, внедрение новой технологии привело к понижению производительности труда.

Сценарии внедрения новой технологии, однако, не ограничиваются ситуацией создания большей массы добавленной стоимости при неизменном количестве занятых на производстве рабочих. Повышение производительности может быть реализовано так, что количество рабочих будет сокращено ($n_2 < n_1$), общее количество добавленной стоимости при этом останется без изменений ($V_2 = V_1$) или их отношение в некоторых рамках будет иметь произвольный характер ($V_2 > V_1$ или $V_2 < V_1$). В таком случае эффект внедрения технологии нельзя рассматривать как простое абсолютное приращение добавленной стоимости, т.~к. оно может оказаться равным нулю ($V_2 = V_1$) даже в том случае, если произошло действительное повышение производительности труда.

Возьмем отношение $V_1 / n_1$ — такое количество добавленной стоимости производит в среднем один рабочий при исходной производительности. В том случае, если бы производительность осталась неизменной, но количество рабочих на производстве изменилось, новое количество рабочих $n_2$ произвело бы добавленной стоимости при старой производительности $(V_1 / n_1) \times n_2$. Вместе с тем, производительность, в действительности, изменилась, $n_2$ рабочих при новой производительности произвели добавленной стоимости $V_2$. Приращение добавленной стоимости, обусловленное измененной производительностью, составит:

$$E = V_2 - (V_1 / n_1) \times n_2$$

Это добавочная добавленная стоимость — прямой экономический эффект $E$ внедрения новой технологии при изменении количества рабочих с $n_1$ на $n_2$. В том случае, если эффект $E$ больше нуля, производительность труда повысилась, если $E = 0$, осталась без изменений, если $E < 0$, внедрение технологии привело к понижению производительности труда.

В том случае, если количество рабочих осталось без изменений ($n_2 = n_1$), экономический эффект, вычисленный по такой формуле, превратится в простое абсолютное приращение добавленной стоимости, рассмотренное выше. В том случае, если количество рабочих сократилось ($n_2 < n_1$), а количество произведенной добавленной стоимости осталось без изменений ($V_2 = V_1$), вычисленный по такой формуле эффект $E$ не обратится в ноль. Формула позволит вычислить эффект $E$ для любых произвольных значений добавленной стоимости $V_1$, $V_2$ и количества рабочих $n_1$, $n_2$, в том числе в случаях, когда количество рабочих увеличилось ($n_2 > n_1$).

Количество добавленной стоимости, производимой рабочим в рассматриваемый период времени, однако, может измениться не только благодаря внедрению новых технологий производства или изменениям интенсивности труда. Сокращение рабочего времени внутри некоторого относительно длительного периода приведет к уменьшению, увеличение рабочего времени приведет к увеличению массы добавленной стоимости, произведенной рабочим. Метрика производительности, основанная на количестве занятых на производстве рабочих, в первом случае покажет понижение производительности труда, во втором — повышение производительности труда, хотя, в действительности, технологии не были внедрены, изменений производительности труда не произошло.

Метод вычисления производительности вычислением произведенного продукта в пересчете на рабочие часы не будет подвержен эффектам растяжения и сжатия рабочего дня, не имеющих прямого отношения к действительной производительности труда рабочего, обусловленной внедренными на производстве технологиями.

В Трудовом кодексе РФ закреплена обязанность организаций вести учет рабочего времени сотрудников (в редакции 2001 года глава 15, статья 91 \cite{tkRf_2001}). Рекомендуемая форма для учета рабочего времени в электронной форме Т-13 «Табель учета рабочего времени» показывает должность сотрудника и общее количество часов, проведенных им на рабочем месте \cite{rosstat_postanov1_05012004_FormT13}. Современные электронные системы учета рабочего времени, внедряемые по собственной инициативе организаций, позволяют учитывать рабочее время сотрудников с еще большей детализацией, в отдельных случаях — до отдельных задач с точностью до минут.

Таким образом, организации обладают всей необходимой информацией для расчета производительности в пересчете на затраты рабочего времени. Производительность труда на предприятии в таком случае может быть определена как отношение произведенной добавленной стоимости к количеству рабочих человеко-часов, затраченных на её производство.

Пусть $V_1$ — добавленная стоимость, произведенная на предприятии за некоторый период до внедрения новой технологии, при этом на производстве для её создания было потрачено $t_1$ рабочих человеко-часов. После внедрения новой технологии за аналогичный период времени была произведена стоимость $V_2$ затратой $t_2$ рабочих человеко-часов.

Если $t_1$ равно $t_2$, т.~е. количество рабочих человеко-часов не изменилось, прямой экономический эффект внедрения технологии составит простое приращение добавленной стоимости:

$$E = V_2 - V_1$$

В этом случае теперь не важно, какое количество рабочих трудилось на предприятии в первом и втором периоде. Стоимость $V_2$ могли произвести то же количество рабочих, что произвели $V_1$, или $V_2$ могли произвести вдвое меньшее количество рабочих за счет переработок в обычные рабочие дни и сверхурочной работы в праздники и выходные, во всех случаях затраченное рабочее время будет $t_2$ человеко-часов. Сокращение или увеличение рабочего дня, таким образом, не влияет на этот метод определения производительности.

Для ситуации, когда после внедрения технологии изменилось общее количество рабочего времени, затрачиваемое на производство продукта на предприятии, т.~е. $t_2$ не равно $t_1$, возьмем сначала отношение $V_1 / t_1$ — такое количество добавленной стоимости будет произведено в среднем за один рабочий человеко-час при исходной производительности.

В том случае, если бы производительность осталась неизменной, но количество примененных рабочих человеко-часов на производстве изменилось, за новое количество рабочих человеко-часов $t_2$ рабочие произвели бы добавленной стоимости при старой производительности $(V_1 / t_1) \times t_2$. Вместе с тем, производительность, в действительности, изменилась, за $t_2$ рабочих человеко-часов при новой производительности рабочие произвели добавленной стоимости $V_2$. Приращение добавленной стоимости, обусловленное измененной производительностью, составит:

$$E = V_2 - (V_1 / t_1) \times t_2$$

Это добавочная добавленная стоимость — прямой экономический эффект $E$ внедрения новой технологии при изменении количества рабочих человеко-часов с $t_1$ на $t_2$. В том случае, если эффект $E$ больше нуля, производительность труда повысилась, если $E = 0$, осталась без изменений, если $E < 0$, внедрение технологии привело к понижению производительности труда.

В том случае, если среднее рабочее время в пересчете на одного рабочего (продолжительность рабочего дня, выходные, праздники и т.~п.) в рассматриваемые периоды времени осталось без изменений, эта формула даст тот же результат, что и предыдущая формула, которая вычисляла экономический эффект на основе производительности в пересчете на количество занятых на производстве рабочих. В случае с новой формулой не важно, какое количество рабочих трудилось на предприятии в первом и втором периоде. К примеру, если технологии производственного процесса оставили без изменений, при этом на втором периоде времени количество рабочих сократили в два раза, а оставшейся половине увеличили рабочее время в два раза за счет увеличения рабочего дня и работы в выходные и праздники, метрика производительности на основе затрат рабочего времени не покажет изменений, отражая действительность, метрика производительности на основе количества рабочих покажет двукратный рост, не соответствующий действительной производительности труда.

Из формулы также видно, что простое наращивание производства без повышения производительности — пропорциональное увеличение затрат рабочей силы и, как следствие, произведенной добавленной стоимости, — даст нулевой экономический эффект. К примеру, если $V_2 = 3V_1$, $t_2 = 3t_1$, $E = 3V_1 - (V_1 / t_1) \times 3t_1 = 0$.

\section*{\raggedright{Сокращение затрат рабочего времени}}

В рассмотренных случаях производительность вычисляется на основе добавленной стоимости, добавленная стоимость по методикам ОЭСР, МОТ, Росстата, Минэкономразвития вычисляется как часть выручки, т.~е. берется в денежной форме. Это значит, что добавленная стоимость, а значит и выводимая из нее метрика производительности, может уменьшиться или увеличиться под действием внешних факторов таких, как колебания цен на производимый товар под действием тех или иных факторов на рынке, не имеющих прямого отношения к производительности предприятия. К примеру, для предприятий, экспортирующих продукт, на объем рублевой выручки будут влиять колебания курса валют.

Для того, чтобы купировать влияние фактора колебания цен на метрику производительности, Минэкономразвития, Росстат, ОЭСР, МОТ применяют специальные поправочные коэффициенты — дефляторы. Дефлятор позволяет привести рыночные цены на продукт в интересующий момент времени к ценам в некоторый предшествующий этому моменту базовый год. Росстат рассчитывает ежегодные значения дефляторов для экономики России в целом и для отдельных отраслей \cite{rosstatGdpDeflatorsData2024}. Однако динамика цены на отдельный товар внутри отрасли или цены на продукт отдельного предприятия, ведущего собственную ценовую политику, может отличаться от усредненной динамики цен внутри всей отрасли, включающей множество наименований товаров. Чем сильнее динамика цены на отдельный интересующий товар отличается от динамики цен внутри отрасли, тем сильнее применение обобщенного дефлятора исказит действительную динамику изменений производительности труда при производстве этого товара внутри отрасли или на конкретном предприятии.

Даже если оставить факторы такого рода в стороне, денежная метрика производительности на основе выручки в общем случае не может отразить изменение производительности на конкретном участке производства, на котором внедрена новая технология. Если предприятие производит несколько видов товаров на разных производственных линиях, можно разделить выручку по этим отдельным видам товаров и оценивать динамику производительности отдельно для каждого из них. Но для одного вида товара оценить изменение производительности с привязкой к отдельным участкам производственного процесса разделением выручки уже нельзя. Если на одном участке внедрена технология, которая повысила производительность, а на другом участке производительность упала по некоторым причинам, общее количество произведенной добавленной стоимости может остаться без изменений, оценить оба этих события по отдельности приведенным способом окажется затруднительно.

Вместе с тем, прямой непосредственный эффект повышения производительности — это сокращение необходимого рабочего времени для производственного процесса. Процесс производства происходит внутри предприятия, произведенный предприятием продукт превращается в деньги за его пределами. Денежная метрика может оценить производительность только после того, как продукт вышел за пределы предприятия. Внутри предприятия производительность труда может быть посчитана на отдельных участках производства в часах.

Добавленные затраты труда, выраженные в рабочих часах, — это время, которое производительные рабочие предприятия тратят на производство целевого продукта, т.~е. это время, которое рабочие проводят на своем рабочем месте, выполняя непосредственные производственные функции, в пересчете на объем произведенного продукта. Чем меньше затрат труда приходится на единицу продукта — единицу количества, веса, длины, площади, объема, энергии и т.~п., — тем выше производительность, чем больше затрат труда, тем ниже производительность. В такой форме затраты могут быть рассчитаны как для готового продукта в финальной точке производственного цикла, так и для любой отдельной промежуточной стадии производства — в пересчете на отдельную составную деталь или на последовательность некоторых трудовых операций.

Информацию об объеме производства продукта — финального или составных элементов или полуфабрикатов промежуточных стадий, — предприятия хранят и накапливают в цифровых системах учета ресурсов предприятия (ERP). Информацию о рабочем времени накапливают в цифровых системах учета рабочего времени. В том случае, если каждый рабочий выполняет производственные функции только на одном участке в общем цикле производства, для оценки общих затрат труда на этих участках хватит данных, собираемых по форме Т-13. В том случае, если некоторые рабочие могут выполнять функции на разных участках производства, потребуется система учета рабочего времени, позволяющая детализировать затраты времени по отдельным задачам, выполняемым конкретными исполнителями.

Пусть $T_1$ — количество рабочих часов, затраты рабочего времени за некоторый период времени до внедрения новой технологии при выполнении некоторой рабочей операции — производстве некоторой детали изделия, выполнении некоторых операций над заготовкой изделия или производстве изделия целиком. $T_2$ — затраты рабочего времени за аналогичный период после внедрения технологии. В первый период времени произведено продукта $m_1$ единиц — штук, кг, метров и т.~п. Во второй период произведено продукта $m_2$ единиц. Величина $T_1 / m_1$ — добавленные затраты рабочего времени в единице продукта при начальной производительности до внедрения новой технологии. $T_2 / m_2$ — добавленные затраты рабочего времени в единице продукта при новой производительности после внедрения новой технологии.

При старой производительности затраты рабочего времени на производство продукта $m_2$ составили бы $(T_1 / m_1) \times m_2$ человеко-часов.

Однако, при новой производительности на продукт $m_2$, в действительности, потрачено $T_2$ человеко-часов. Таким образом, при изменении объема производства продукта с $m_1$ до $m_2$ экономический эффект от внедрения новой технологии, выраженный в экономии труда, составит:

$$E = (T_1 / m_1) \times m_2 - T_2 \quad\mbox{человеко-часов}$$

В том случае, если эффект внедрения $E$ больше нуля, производительность труда повысилась, в том случае, если $E$ равен нулю, производительность осталась без изменений, если $E$ меньше нуля, внедрение новой технологии привело к понижению производительности труда.

В том случае, если объем производства продукта не изменился ($m_2 = m_1$), экономический эффект составит простое сокращение затрат труда ($T_1 - T_2$) в известном объеме продукта $m_1$.

В том случае, если технологический процесс претерпевает комплексные преобразования, одна или несколько технологий внедрены на разных участках производства, общий суммарный экономический эффект составит:

$$E_{\mbox{сум}} = E^{(1)} + … + E^{(n)}$$

Где $E^{(i)}$ — эффект на $i$-м участке производства, $i = 1,…,n$, $n$ — количество точек внедрения.

Повышение производительности труда на отдельном участке производства, тем не менее, не означает автоматическое повышение производительности при выпуске конечного изделия. К примеру, если реализовать повышение производительности наращиванием производства отдельных деталей на соответствующем участке производства (увеличить $m_2$ на $i$-м участке), новые добавочные детали могут оказаться невостребованными при производстве финальной массы конечного продукта, если для создания добавочной массы конечного продукта не хватит новых добавочных деталей из других точек производственного процесса.

Экономический эффект при производстве конечного продукта $E_{\mbox{кп}}$, рассчитанный по формуле экономии труда, в общем случае не обязательно будет равен суммарному экономическому эффекту $E_{\mbox{сум}}$, полученному суммой эффектов во всех промежуточных точках внедрения новых технологий. Затраты рабочего времени $T_2$ и масса произведенного продукта $m_2$ на каждой $i$-й стадии производства после внедрения новой технологии должны быть подобраны таким образом, чтобы сумма всех экономических эффектов во всех промежуточных точках внедрения совпала с экономическим эффектом при производстве новой массы $m_2$ конечного продукта.

Пусть изделие — конечный продукт, — состоит из 2-х деталей, \textit{деталь-1} делают 6 рабочих за 8 часов ($m_1^{(1)} = 1$ шт, $T_1^{(1)}=48$ человеко-часов), \textit{деталь-2} делает 1 рабочий за 8 часов ($m_1^{(2)} = 1$ шт, $T_1^{(2)}=8$ человеко-часов), из двух деталей собирают изделие еще 1 рабочий еще за 8 часов ($m_1^{(3)} = 1$ шт, $T_1^{(3)}=8$ человеко-часов). Производственный процесс организован таким образом, что на последнем участке сборки в начале каждого нового дня всегда есть готовая \textit{деталь-1} и готовая \textit{деталь-2}, произведенные днем ранее, поэтому каждый новый день 2 рабочих финальной стадии выпускают 1 конечное изделие, собранное из 2-х деталей, произведенных в предыдущий день. За четыре пятидневные рабочие недели (всего 20 рабочих дней) с восьмичасовым рабочим днем (всего 160 часов для одного человека) в модуле складского учета системы ERP будет показано количество произведенных \textit{деталей-1} — $m_1^{(1)} = 20$ штук, количество произведенных \textit{деталей-2} — $m_1^{(2)} = 20$ штук, количество произведенных конечных изделий — $m_1 = m_1^{(3)} = 20$ штук (в первый день цикла на складе доступны одна \textit{деталь-1} и \textit{деталь-2}, которые можно задействовать при производстве конечного изделия в этот же день, в последний день цикла произведенные \textit{деталь-1} и \textit{деталь-2} не задействованы при производстве последней единицы конечного продукта, остаются на складе на следующий цикл). Система учета рабочего времени отразит затраты труда на отдельных участках производственной линии: при производстве \textit{детали-1} $T_1^{(1)} = 160 \times 6 = 960$ человеко-часов, при производстве \textit{детали-2} $T_1^{(2)} = 160 \times 1 = 160$ человеко-часов, на этапе сборки конечного изделия $T_1^{(3)} = 160 \times 1 = 160$ человеко-часов, всего $T_1 = 960 + 160 + 160 = 1280$ человеко-часов.

Добавленные затраты рабочего времени в единице конечного продукта в таком случае составят: $T_1 / m_1 = 1280 / 20 = 64$ человеко-часа.

Пусть теперь в производственный процесс внедрена некоторая технология, которая повысила производительность труда на участке производства \textit{детали-1} в 3 раза, на остальных участках производительность осталась без изменений. Теперь на производство одной \textit{детали-1} требуется затратить $48 / 3 = 16$ человеко-часов. Допустим, в условиях новой производительности, предприятие решило оставить на участке производства \textit{детали-1} прежнее количество рабочих с прежней продолжительностью рабочего времени. В конце нового производственного цикла в модуле складского учета системы ERP будет показано количество произведенных \textit{деталей-1} — $m_2^{(1)} = 60$ штук (новая производительность), количество произведенных \textit{деталей-2} — $m_2^{(2)} = 20$ штук (производительность не изменилась). Для производства конечного изделия всё еще требуется одна \textit{деталь-1} и одна \textit{деталь-2}, поэтому максимальное количество произведенных конечных изделий ограничено, во-первых, наименьшим из доступных наборов составных деталей, в этом случае — количеством доступных \textit{деталей-2}, во-вторых, производительностью и количеством рабочих на соответствующем участке сборки: $m_2 = m_2^{(3)} = 20$ штук (производительность не изменилась).

Система учета рабочего времени отразит затраты труда на отдельных участках производственной линии: при производстве \textit{детали-1} $T_2^{(1)} = 160 \times 6 = 960$ человеко-часов, при производстве \textit{детали-2} $T_2^{(2)} = 160 \times 1 = 160$ человеко-часов, на этапе сборки конечного изделия $T_2^{(3)} = 160 \times 1 = 160$ человеко-часов, всего $T_2 = 960 + 160 + 160 = 1280$ человеко-часов.

Добавленные затраты рабочего времени в единице конечного продукта в таком случае составят: $T_2 / m_2 = 1280 / 20 = 64$ человеко-часа, — столько же, сколько было до внедрения новой технологии. Полный экономический эффект в таком случае, вполне очевидно, окажется равен нулю:

$$E = (T_1 / m_1) \times m_2 - T_2 = (1280 / 20) \times 20 - 1280 = 0 \quad\mbox{человеко-часов}$$

Таким образом, не смотря на то, что в процесс производства была внедрена новая технология, и производительность труда на одном из участков, действительно, повысилась, при выбранной стратегии внедрения затраты труда в единице конечного продукта, рассчитанные по такой методике, остались прежними.

Рассмотрим добавленные затраты труда отдельно на каждом из участков производства. В первый исследуемый цикл производства до внедрения технологии добавленные затраты труда составят: на одну \textit{деталь-1} — $T_1^{(1)} / m_1^{(1)} = 960 / 20 = 48$ человеко-часов, на одну \textit{деталь-2} — $T_1^{(2)} / m_1^{(2)} = 160 / 20 = 8$ человеко-часов, на этап сборки одного конечного изделия — $T_1^{(3)} / m_1^{(3)} = 160 / 20 = 8$ человеко-часов, суммарные добавленные затраты труда в одном конечном изделии — $48 + 8 + 8 = 64$ человеко-часа.

Во второй исследуемый цикл производства после внедрения технологии добавленные затраты труда составят: на одну \textit{деталь-1} — $T_2^{(1)} / m_2^{(1)} = 960 / 60 = 16$ человеко-часов, на одну \textit{деталь-2} — $T_2^{(2)} / m_2^{(2)} = 160 / 20 = 8$ человеко-часов, на этап сборки одного конечного изделия — $T_2^{(3)} / m_2^{(3)} = 160 / 20 = 8$ человеко-часов, суммарные добавленные затраты труда в одном конечном изделии — $16 + 8 + 8 = 32$ человеко-часа.

Экономический эффект на первом участке при производстве \textit{детали-1}:

\begin{align}
E^{(1)} &= (T_1^{(1)} / m_1^{(1)}) \times m_2^{(1)} - T_2^{(1)} =\nonumber\\
        &= (960 / 20) \times 60 - 960 = 1920 \quad\mbox{человеко-часов}\nonumber
\end{align}

На втором участке при производстве \textit{детали-2}:

\begin{align}
E^{(2)} &= (T_1^{(2)} / m_1^{(2)}) \times m_2^{(2)} - T_2^{(2)} =\nonumber\\
        &= (160 / 20) \times 20 - 160 = 0 \quad\mbox{человеко-часов}\nonumber
\end{align}

На последнем участке при сборке конечного изделия:

\begin{align}
E^{(3)} &= (T_1^{(3)} / m_1^{(3)}) \times m_2^{(3)} - T_2^{(3)} =\nonumber\\
        &= (160 / 20) \times 20 - 160 = 0 \quad\mbox{человеко-часов}\nonumber
\end{align}

Суммарный экономический эффект на всех участках:

$$E_{\mbox{сум}} = E^{(1)} + E^{(2)} + E^{(3)} = 1920 + 0 + 0 = 1920 \quad\mbox{человеко-часов}$$

Таким образом, как добавочные затраты труда в пересчете на одну деталь, так и общий экономический эффект от внедрения технологии, будут отличаться в зависимости от того, какой способ их вычисления выбран. Метод вычисления производительности, берущий за основу производительность на отдельных участках производства, показывает экономический эффект — 1920 часов сэкономленного труда. В то же время метод вычисления производительности, берущий за основу общие затраты труда, показывает нулевой экономический эффект, сэкономленный труд не показывает. Легко видеть, что при выбранном сценарии внедрения технологии, сэкономленный труд, который видно при вычислении производительности на отдельных участках, воплотился в производстве добавочного количества \textit{деталей-1} (40 штук), которые, в итоге, так и не были задействованы при производстве конечных изделий.

С одной стороны, может показаться, что метод, учитывающий производительность отдельных участков, является более точным, т.~к. он отразил действительное повышение производительности и сэкономленный труд внутри конечного изделия. С другой стороны, если не брать в рассмотрение ситуацию, когда произведенная добавочная масса \textit{деталей-1} является отдельным самодостаточным товаром, без перестройки производства вся эта произведенная добавочная масса промежуточных деталей обернется потерями. В том случае, если произведенные детали — это какие-то относительно долговечные изделия, например, металлические заготовки, они будут заполнять склад, но, рано или поздно, все равно могут подвергнуться коррозии и прочей деградации, в любом случае будут лежать без применения. В том случае, если это некоторый скоропортящийся продукт, или, например, избыточный выпуск тепла, электричества или прочей энергии, потери при их сиюминутном неприменении по назначению становятся еще более очевидными. Труд, потраченный на их производство, пропадет и, будучи частью единого производственного процесса внутри предприятия, всё равно войдет в стоимость конечного изделия. Итоговый экономический эффект обратится в ноль, как и показывает метод вычисления производительности, берущий за основу общие итоговые затраты.

Хороший сценарий внедрения технологии приведет к тому, что промежуточные локальные экономические эффекты на отдельных участках производства проявят себя в конечном общем экономическом эффекте производства конечного изделия. В идеальном случае суммарный экономический эффект, вычисленный как сумма локальных эффектов, и общий экономический эффект должны совпасть. Это произойдет в том случае, если после внедрения новой технологии производственные процессы на целевом участке все ещё будут согласованы с прочими участками производства, оставшимися без изменений, или же производственные процессы будут перестроены на всей производственной линии так, что их протекание будет снова согласовано между собой и при этом будет учитывать новую повышенную производительность отдельных точек производства.

Рассмотрим другой сценарий внедрения новой технологии, которая, как и в рассмотренном случае, повышает производительность труда на участке производства \textit{детали-1} в 3 раза (на производство одной \textit{детали-1} требуется затратить $48 / 3 = 16$ человеко-часов). На этот раз предприятие решило сократить количество задействованного на участке производства \textit{детали-1} труда, сохранив при этом прежний выпуск этих деталей. 

В конце нового производственного цикла в модуле складского учета системы ERP будет показано количество произведенных \textit{деталей-1} — $m_2^{(1)} = 20$ штук (новая производительность), количество произведенных \textit{деталей-2} — $m_2^{(2)} = 20$ штук (производительность не изменилась), количество произведенных конечных изделий — $m_2 = m_2^{(3)} = 20$ штук (производительность не изменилась). Система учета рабочего времени отразит затраты труда на отдельных участках производственной линии: при производстве \textit{детали-1} $T_2^{(1)} = 16$ часов $\times 20$ дней $= 320$ человеко-часов (одна деталь в день, 20 рабочих дней), при производстве \textit{детали-2} $T_2^{(2)} = 160 \times 1 = 160$ человеко-часов, на этапе сборки конечного изделия $T_2^{(3)} = 160 \times 1 = 160$ человеко-часов, всего $T_2 = 320 + 160 + 160 = 640$ человеко-часов.

Добавленные затраты рабочего времени в единице конечного изделия в таком случае составят: $T_2 / m_2 = 640 / 20 = 32$ человеко-часа. Полный экономический эффект:

$$E = (T_1 / m_1) \times m_2 - T_2 = (1280 / 20) \times 20 - 640 = 640 \quad\mbox{человеко-часов}$$

Рассмотрим теперь добавленные затраты труда на каждом отдельном участке производства. Во второй исследуемый цикл производства после внедрения технологии добавленные затраты труда составят: на одну \textit{деталь-1} — $T_2^{(1)} / m_2^{(1)} = 320 / 20 = 16$ человеко-часов, на одну \textit{деталь-2} — $T_2^{(2)} / m_2^{(2)} = 160 / 20 = 8$ человеко-часов, на этап сборки одного конечного изделия — $T_2^{(3)} / m_2^{(3)} = 160 / 20 = 8$ человеко-часов, суммарные добавленные затраты труда в одном конечном изделии — $16 + 8 + 8 = 32$ человеко-часа.

Экономический эффект на первом участке при производстве \textit{детали-1}:

\begin{align}
E^{(1)} &= (T_1^{(1)} / m_1^{(1)}) \times m_2^{(1)} - T_2^{(1)} =\nonumber\\
        &= (960 / 20) \times 20 - 320 = 640 \quad\mbox{человеко-часов}\nonumber
\end{align}

На втором участке при производстве \textit{детали-2}:

\begin{align}
E^{(2)} &= (T_1^{(2)} / m_1^{(2)}) \times m_2^{(2)} - T_2^{(2)} =\nonumber\\
        &= (160 / 20) \times 20 - 160 = 0 \quad\mbox{человеко-часов}\nonumber
\end{align}

На последнем участке при сборке конечного изделия:

\begin{align}
E^{(3)} &= (T_1^{(3)} / m_1^{(3)}) \times m_2^{(3)} - T_2^{(3)} =\nonumber\\
        &= (160 / 20) \times 20 - 160 = 0 \quad\mbox{человеко-часов}\nonumber
\end{align}

Суммарный экономический эффект на всех участках:

$$E_{\mbox{сум}} = E^{(1)} + E^{(2)} + E^{(3)} = 640 + 0 + 0 = 640 \quad\mbox{человеко-часов}$$

Таким образом, экономический эффект, рассчитанный на основе общих затрат труда, на этот раз получился не нулевым и совпал с экономическим эффектом, рассчитанным как сумма локальных экономических эффектов на отдельных участках производства. Внедрение технологии на отдельном участке производства повлияло на общую производительность производства конечного продукта, т.~к. технология была внедрена таким способом, что все участки производства оказались снова согласованы — для производства конечного продукта на последнем этапе были использованы все имеющиеся в наличии \textit{детали-1} и \textit{детали-2}.

Зная особенности технологического процесса, осуществить согласование между отдельными его участками можно подгоняя объемы производства на отдельных связанных между собой промежуточных участках. В частности, подбирая значения параметров модели — количество затрат труда на каждом из рассматриваемых участков $T^{(i)}$ и количество произведенного на каждом из этих участков промежуточного продукта $m^{(i)}$. В том случае, если стоит задача оценить эффективность постфактум после состоявшегося внедрения технологии, конкретные значения количества затрат труда и произведенного продукта можно брать как статистическую данность из наблюдений за целевым участком производства. В том случае, если стоит задача оценить и сравнить разные сценарии до непосредственного внедрения технологии, их следует брать как целевые и расчетные параметры. Для каждого уровня производительности, обусловленного той или иной технологией, между затратами труда и произведенным их применением объемом продукта будет установлена некоторая функциональная связь. При данной производительности известное количество затрат труда $T^{(i)}$ воплотится в объеме продукта $m^{(i)}$, таким образом, между $m^{(i)}$ и $T^{(i)}$ можно определить прямую функциональную зависимость: $m = m^{(i)}(T)$. Аналогичным образом, для производства известного объема продукта $m^{(i)}$ при данной производительности потребуется количество затрат труда $T^{(i)}$, обратная зависимость будет выглядеть как $T = T^{(i)}(m)$.

В том случае, если при подгонке модели есть необходимость отталкиваться от объемов производства (к примеру, речь об ограниченном заказе при относительно свободном доступе к рабочей силе), за данный целевой параметр следует брать целевой объем продукта $m^{(i)}$, затраты труда будут получены как расчетное значение $T^{(i)}(m^{(i)})$. В том случае, если при подгонке модели есть необходимость отталкиваться от затрат труда (к примеру, количество специалистов на предприятии и на рынке труда ограничено, стоит задача произвести с имеющейся в наличии рабочей силой максимальный или оптимальный объем продукта), за данный целевой параметр следует брать целевые затраты труда на одном или нескольких участках $T^{(i)}$, объем продукта будет получен как расчетное значение $m^{(i)}(T^{(i)})$. Или же варьировать можно и те и другие значения, сохраняя между ними отношение установленной зависимости, чтобы получить некоторый промежуточный оптимальный компромисс.

В первом рассмотренном сценарии внедрения технологии мы отталкивались от задействованного на производстве текущего количества рабочей силы, новый объем производства на целевом промежуточном участке был получен как значение функциональной зависимости. Но после внесения изменений в целевой участок производства с наращиванием производства, не провели согласование объемов производства на всех участках производственной линии до финала, в итоге производственная линия оказалась рассогласована. Во втором сценарии отталкивались от объема продукта, сократив при этом затраты труда. Объемы производства на всех участках линии оказались сразу согласованы, т.~к. за целевой объем промежуточного продукта был взят его текущий выпуск.

Рассмотрим еще один сценарий внедрения новой технологии. На этот раз попробуем сохранить общий объем применяемых на предприятии затрат труда, т.~е. оставить без изменений количество занятых рабочих и продолжительность их рабочего дня, при этом преобразовать производство так, чтобы объем выпуска на всех участках линии оказался вновь согласован уже с учетом новой производительности.

Как и раньше, технология повышает производительность труда на участке производства \textit{детали-1} в 3 раза (на производство одной \textit{детали-1} требуется затратить $48 / 3 = 16$ человеко-часов). Пусть на \textit{участке-1} из 6-ти рабочих останется 4, за полный рабочий день с новой производительностью они произведут 2 \textit{детали-1}. Одного из 2-х освободившихся рабочих отправим на участок производства \textit{детали-2}, вместе с рабочим, который там уже находился, они смогут произвести за рабочий день две \textit{детали-2}. Еще одного освободившегося на \textit{участке-1} рабочего отправим на \textit{участок-3}, вместе с тем рабочим, который там уже находился, они смогут производить в день два конечных изделия из имеющихся теперь в распоряжении двух \textit{деталей-1} и двух \textit{деталей-2}.

В конце нового производственного цикла в модуле складского учета системы ERP будет показано количество произведенных \textit{деталей-1} — $m_2^{(1)} = 40$ штук (новая производительность), количество произведенных \textit{деталей-2} — $m_2^{(2)} = 40$ штук (производительность не изменилась, производство расширено), количество произведенных конечных изделий — $m_2 = m_2^{(3)} = 40$ штук (производительность не изменилась, производство расширено). Система учета рабочего времени отразит затраты труда на отдельных участках производственной линии: при производстве \textit{детали-1} $T_2^{(1)} = 160 \times 4 = 640$ человеко-часов, при производстве \textit{детали-2} $T_2^{(2)} = 160 \times 2 = 320$ человеко-часов, на этапе сборки конечного изделия $T_2^{(3)} = 160 \times 2 = 320$ человеко-часов, всего $T_2 = 640 + 320 + 320 = 1280$ человеко-часов.

Добавленные затраты рабочего времени в единице конечного изделия в таком случае составят: $T_2 / m_2 = 1280 / 40 = 32$ человеко-часа. Полный экономический эффект:

$$E = (T_1 / m_1) \times m_2 - T_2 = (1280 / 20) \times 40 - 1280 = 1280 \quad\mbox{человеко-часов}$$

Рассмотрим теперь добавленные затраты труда на каждом отдельном участке производства. Во второй исследуемый цикл производства после внедрения технологии добавленные затраты труда составят: на одну \textit{деталь-1} — $T_2^{(1)} / m_2^{(1)} = 640 / 40 = 16$ человеко-часов, на одну \textit{деталь-2} — $T_2^{(2)} / m_2^{(2)} = 320 / 40 = 8$ человеко-часов, на этап сборки одного конечного изделия — $T_2^{(3)} / m_2^{(3)} = 320 / 40 = 8$ человеко-часов, суммарные добавленные затраты труда в одном конечном изделии — $16 + 8 + 8 = 32$ человеко-часа.

Экономический эффект на первом участке при производстве \textit{детали-1}:

\begin{align}
E^{(1)} &= (T_1^{(1)} / m_1^{(1)}) \times m_2^{(1)} - T_2^{(1)} =\nonumber\\
        &= (960 / 20) \times 40 - 640 = 1280 \quad\mbox{человеко-часов}\nonumber
\end{align}

На втором участке при производстве \textit{детали-2}:

\begin{align}
E^{(2)} &= (T_1^{(2)} / m_1^{(2)}) \times m_2^{(2)} - T_2^{(2)} =\nonumber\\
        &= (160 / 20) \times 40 - 320 = 0 \quad\mbox{человеко-часов}\nonumber
\end{align}

На последнем участке при сборке конечного изделия:

\begin{align}
E^{(3)} &= (T_1^{(3)} / m_1^{(3)}) \times m_2^{(3)} - T_2^{(3)} =\nonumber\\
        &= (160 / 20) \times 40 - 320 = 0 \quad\mbox{человеко-часов}\nonumber
\end{align}

Суммарный экономический эффект на всех участках:

$$E_{\mbox{сум}} = E^{(1)} + E^{(2)} + E^{(3)} = 1280 + 0 + 0 = 1280 \quad\mbox{человеко-часов}$$

Экономический эффект, рассчитанный на основе общих затрат труда, снова совпал с экономическим эффектом, рассчитанным как сумма локальных экономических эффектов на отдельных участках производства. Хотя внедрение технологии было произведено с увеличением объема производства на целевом участке, производство на всех прочих участках было согласовано с новым объемом. В итоге локальный экономический эффект произвел вклад в повышение общей производительности при производстве конечного изделия, в отличие от первого рассмотренного сценария, в котором вместе с расширением производства на целевом участке производство на прочих участках не было согласовано.

Сценарий внедрения технологии с сохранением прежнего объема производства оказывается проще с точки зрения согласования всех участков производства при новой производительности, т. к. фактически не требует перестройки производства за пределами целевого участка внедрения. Однако, интересно отметить, что внедрение технологии с расширением производства при сохранении прежнего количества затрат труда произвело заметно больший экономический эффект (1280 человеко-часов), чем внедрение этой же технологии с сохранением прежнего объема производства при сокращении затрат труда (640 человеко-часов).

Дело в том, что в сценарии сокращения затрат труда экономический эффект представляет собой добавочный продукт, который после повышения производительности создает оставшаяся на предприятии рабочая сила, и который замещает продукт, создаваемый ранее сокращенной рабочей силой при более низкой производительности, пока она еще не была сокращена. При этом высвободившаяся после повышения производительности рабочая сила в потенциале может быть применена для создания еще большего объема продукта, но в сценарии сокращения это произойдет за пределами предприятия, и поэтому эта доля эффекта не попадет в распоряжение предприятия и в его внутренние системы учета. В сценарии увеличения производства с сохранением общего объема примененных затрат труда эта доля добавочного продукта производится внутри предприятия и составляет часть общего реализованного экономического эффекта. Таким образом, экономический эффект, попадающий в распоряжение предприятия после повышения производительности, в сценарии увеличения объема производства с сохранением затрат труда превышает в два раза эффект в сценарии сокращения затрат труда с сохранением объема производства.

\section*{\raggedright{Превращение сокращения затрат рабочего времени на производстве в деньги}}

Пусть предприятие до внедрения новой технологии работало со средней по отрасли производительностью, т.~е. затраты труда на производство одного изделия на рассматриваемом предприятии были в среднем примерно такими же, как на большинстве других конкурирующих предприятий, производящих идентичный продукт. Допустим, при такой производительности один рабочий за рассматриваемый период (один месяц — 20 рабочих дней) создает в денежном выражении добавленной стоимости 800 тыс. руб. В месяц 8 рабочих создают 20 конечных изделий (шестеро на 1-м участке, по одному на 2-м и 3-м). Добавленной стоимости в месячном конечном продукте в деньгах будет 800 тыс. × 8 = 6400 тыс. руб. (6 млн. 400 тыс. руб.).

Во втором сценарии технология была внедрена таким образом, что на первом участке производства был сохранен выпуск продукта $m_2^{(1)} = m_1^{(1)} = 20$ \textit{деталей-1}, при этом были сокращены затраты труда. Сокращение затрат труда при сохранении объема выпуска продукта в этом сценарии на практике может быть реализовано несколькими способами.

К примеру, из шести изначально задействованных на 1-м участке рабочих предприятие может оставить двоих, остальных четверых сократить. Затраты труда на участке производства \textit{детали-1} за цикл в таком случае будут: $T_2^{(1)}$ = 160 часов × 2 человека = 320 человеко-часов. Освободилось 160 часов × 4 человека = 640 человеко-часов.

Владельцы предприятия могут конвертировать эти сэкономленные часы в добавочную денежную прибыль. При новой производительности месячный продукт предприятия 20 конечных изделий создается трудом 4-х рабочих (двое на 1-м участке, по одному на 2-м и 3-м). Добавленная стоимость, вычисленная на основе затрат труда при новой производительности предприятия, в месячном продукте в деньгах составит 800 тыс × 4 = 3200 тыс. руб. (3 млн. 200 тыс. руб.).

Таким образом, абсолютная величина добавленной стоимости в деньгах, на первый взгляд, уменьшилась с 6400 до 3200 тыс. руб., в соответствии с сокращением задействованной на производстве рабочей силы. Однако, после повышения индивидуальной производительности труда на одном рассматриваемом предприятии, рыночные цены на идентичные продукты других предприятий-конкурентов, обусловленные средними по отрасли затратами труда, останутся прежними. Предприятие всё еще сможет продавать выпускаемый с новой производительностью продукт по старой цене. Значит при продаже месячного продукта 20 изделий в его распоряжении окажутся всё те же 6400 тыс. руб., а не 3200 тыс. руб. добавленной стоимости, но при этом на производстве для создания этого продукта было задействовано не 8 рабочих, а 4 рабочих. Добавочная добавленная стоимость в денежном выражении в распоряжении предприятия в таком случае составит: 6400 тыс. руб - 3200 тыс. руб. = 3200 тыс. руб. Это добавочные денежные средства, которые предприятие извлекает с рынка «просто так» благодаря тому, что индивидуальная производительность предприятия оказалась выше производительности средней по отрасли, они же составляют денежное выражение реализованного внедрением технологии экономического эффекта.

Таким образом, мы разными путями получили экономический эффект в разных формах: в затратах труда 640 человеко-часов как экономию труда сокращенных 4-х рабочих, в денежной форме 3200 тыс. руб. как добавочную добавленную стоимость, извлеченную с рынка благодаря тому, что индивидуальная производительность предприятия опередила производительность предприятий-конкурентов. До внедрения технологии и повышения производительности за один человеко-час один рабочий производил добавленной стоимости в денежном выражении: 800 тыс. руб. / 160 часов = 5 тыс. руб. Четверо тогда еще не сокращенных рабочих создавали добавленной стоимости в месяц общими затратами труда 640 человеко-часов: 5 тыс. руб. × 640 человеко-часов = 3200 тыс. руб. После внедрения технологии эти же денежные средства «создаются» трудом оставшихся на производстве рабочих в дополнение к тем средствам, которые они уже создавали с прежней производительностью, при условии, что рыночная цена на продукт не изменилась. Значит эта денежная величина и составляет экономический эффект внедрения технологии, численно равный результату, полученному выше, но на этот раз полученный из величины экономического эффекта в форме экономии затрат труда.

Таким образом, экономический эффект в денежном выражении можно получить из экономического эффекта в форме экономии рабочего времени по следующей формуле:

$$E_v = v \times E_t$$

Где $v$ — количество добавленной стоимости в денежном выражении, создаваемой рабочим за один человеко-час при средней по отрасли производительности или при исходной производительности предприятия до внедрения новой технологии. $E_t$ — количество человеко-часов, сэкономленных благодаря внедрению технологии, или экономический эффект в форме экономии затрат труда.

Допустим, зарплата одного рабочего за рассматриваемый период времени (месяц — 20 рабочих дней) составляет 200 тыс. руб. Фонд оплаты труда (ФОТ) до внедрения технологии на предприятии с 8-ю рабочими составлял 200 тыс. × 8 = 1600 тыс. руб. После внедрения технологии и сокращения 4-х рабочих на предприятии осталось всего 4 рабочих, в ФОТ будет направлено 200 тыс. руб. × 4 = 800 тыс. руб. Таким образом, сокращение ФОТ после внедрения технологии составит 1600 тыс. руб. - 800 тыс. руб. = 800 тыс. руб. При том, что выручка (как и входящая в нее добавленная стоимость), извлекаемая с рынка, при неизменном выпуске продукта и неизменных ценах осталась прежней, эти 800 тыс. руб., полученные от сокращения ФОТ, предстанут перед владельцами предприятия как вся новая добавочная прибыль — новые свободные деньги, которые после внедрения технологии и сокращения персонала попали в их распоряжение. Однако, легко видеть, что сумма полного экономического эффекта 3200 тыс. руб., вычисленная ранее, значительно превышает простую экономию денежных средств от сокращения ФОТ 800 тыс. руб. Таким образом, действительный экономический эффект от внедрения технологии в сценарии сокращения затрат труда не ограничивается простым сокращением фонда оплаты труда сокращенной рабочей силы.

Такой результат может показаться не вполне интуитивным. Для лучшей наглядности можно рассмотреть аналогичную ситуацию, но с другими начальными условиями. Пусть на соседнем предприятии с полностью идентичным производственным процессом до внедрения технологии изначально трудилось 4 рабочих (столько же, сколько на рассматриваемом предприятии осталось после внедрения технологии), производимая их трудом добавленная стоимость в деньгах составляла 800 тыс. руб. × 4 = 3200 тыс. руб. После внедрения этой же технологии это предприятие решило не сокращать персонал, а увеличить выпуск продукта при старом количестве рабочих. Таким образом, в результате внедрения одной и той же технологии изначально рассматриваемое и «соседнее» предприятия оказались полностью идентичны — и по производительности труда, и по объему производимого продукта, и по количеству задействованных на производстве человек. Количество добавленной стоимости, извлекаемой обоими предприятиями с рынка, в денежном выражении составит 6400 тыс. руб. Для «соседнего» предприятия количество добавочной добавленной стоимости составит 6400 тыс. руб. - 3200 тыс. руб. = 3200 тыс. руб., для него это будет экономический эффект от внедрения новой технологии в денежном выражении. Он оказался в точности таким же, какой был достигнут на исходном предприятии сокращением рабочих. При этом легко видеть, что для «соседнего» предприятия экономический эффект никак не зависит от ФОТ, который вообще не был сокращен.

Таким образом, капитал, вложенный в предприятие с повышенной производительностью, работает более эффективно. Экономический эффект отражает эту эффективность в сравнении с предприятием — этим же до модернизации или другим, — работающим с меньшей производительностью.

Другой вариант реализации сокращения затрат труда во втором сценарии внедрения технологии — сохранить исходное количество рабочих 6 человек на 1-м участке производства, но при этом сократить их рабочий день пропорционально повышению производительности труда — в 3 раза (рабочий день одного человека: 8 часов / 3 = 2.7 часа, 20 рабочих дней — 160 часов / 3 = 53.3 часа). Затраты труда на 1-м участке производства за цикл будут: $T_2^{(1)}$ = (160 часов / 3) × 6 человек = 320 человеко-часов. Освободилось \{960 человеко-часов (прошлые затраты) - 320 человеко-часов (новые затраты) =\} 640 человеко-часов. На 2-м и 3-м участке работает по одному рабочему, их рабочий день остается без изменений, экономии труда на этих участках нет, общая экономия труда при производстве конечного изделия — 640 человеко-часов.

При новой производительности месячный продукт предприятия 20 конечных изделий создается трудом 8-ми рабочих, 6 из которых работают на \textit{участке-1} $1/3$ восьмичасового рабочего дня, по одному работают на 2-м и 3-м участках полный восьмичасовой рабочий день. Если один рабочий за полный восьмичасовой рабочий день создает в денежном выражении 800 тыс. руб. добавленной стоимости, за $1/3$ рабочего дня он создаст 800 тыс. руб. / 3 добавленной стоимости. Добавленная стоимость, вычисленная на основе затрат труда при новой производительности предприятия, в месячном продукте на \textit{участке-1} в деньгах составит (800 тыс. руб. / 3) × 6 = 1600 тыс. руб., в конечном месячном продукте в деньгах составит \{1600 тыс. руб. (\textit{участок-1}) + 800 тыс. руб. (\textit{участок-2}) + 800 тыс. руб. (\textit{участок-3}) =\} 3200 тыс. руб. Вполне очевидно, что она равна добавленной стоимости, созданной 4-мя рабочими, работающими полный рабочий день, как и в предыдущем варианте реализации сценария сокращения затрат труда.

Величина добавленной стоимости в выручке от продажи конечного продукта 20 изделий при прежней цене на рынке всё также составит 6400 тыс. руб. Экономический эффект, таким образом, составит 6400 тыс. руб. - 3200 тыс. руб. = 3200 тыс. руб., как и в предыдущем варианте реализации сценария. Этот же результат можно получить из представленной выше формулы:

$$E_v = v \times E_t = 5 \quad\mbox{тыс. руб.} \times 640 \quad\mbox{человеко-часов} = 3200 \quad\mbox{тыс. руб.}$$

В том случае, если сокращение рабочего дня сопровождается пропорциональным сокращением заработной платы, владельцы предприятия получат в свое распоряжение весь экономический эффект, реализованный внутри предприятия, целиком. Перешедшие на неполную ставку рабочие при этом будут вынуждены применить «освобожденные» часы на другом рабочем месте за пределами предприятия (на другом производстве или в других отраслях экономики), чтобы компенсировать потерю доли заработка и сохранить возможность получения необходимых жизненных средств.

Размер заработной платы одного рабочего на \textit{участке-1} в месяц при неполном рабочем дне составит 200 тыс. руб. / 3 = 66.667 тыс. руб. Размер ФОТ для 6-ти рабочих при полном рабочем дне составит 200 тыс. руб. × 6 = 1200 тыс. руб. Величина потенциала сокращения ФОТ при переходе на неполный рабочий день для 6-ти рабочих на \textit{участке-1} составит 1200 тыс. руб. - (200 тыс. руб. / 3) × 6 = 800 тыс. руб.

В том случае, если сокращение рабочего дня произойдет с сохранением прежнего уровня заработной платы, доля экономического эффекта, реализованного внутри предприятия, перейдет в руки рабочих \textit{участка-1} в размере потенциала сокращения ФОТ 800 тыс. руб., ФОТ при этом не будет сокращен. В руки владельцев предприятия при этом всё еще попадет доля экономического эффекта сверх потенциала сокращения ФОТ 3200 тыс. руб. - 800 тыс. руб. = 2400 тыс. руб. Кроме того, перешедшие на неполную ставку рабочие, сохранившие при этом прежний денежный заработок, обеспечивающий им получение всех необходимых жизненных средств, не будут вынуждены искать дополнительное применение своей рабочей силы за пределами предприятия. Таким образом, в их распоряжение перейдет доля эффекта сокращения затрат труда в первоначальной форме — добавочного количества свободного времени, которое они могут тратить по своему усмотрению. Это время может быть учтено за пределами предприятия в масштабах экономики как пополнение фонда общественного времени досуга (всего 640 часов на рассмотренном периоде). 

Можно также преподнести сокращение затрат труда так, что количество рабочих останется без изменений, продолжительность рабочего дня останется тоже без изменений, но для того, чтобы избежать перепроизводства неиспользуемых промежуточных деталей, их производство можно организовать прерывающимися периодами внутри полного цикла: в течение одной части полного цикла рабочие на \textit{участке-1} производят детали с повышенной производительностью, избыток деталей накапливается на складе, в следующей части цикла на \textit{участке-1} делается перерыв, рабочие остальных участков производства потребляют запас произведенных избыточных деталей, далее рабочие \textit{участка-1} снова выходят на производство на полный рабочий день, производят новый запас \textit{деталей-1} и т.~п. В действительности, этот вариант полностью идентичен предыдущему варианту с переводом рабочих на частичную занятость, только с другим графиком работы, т.~к. на полном цикле будет потрачено то же самое соответствующее новой производительности сокращенное количество труда.

В третьем сценарии технология была внедрена таким образом, что количество задействованной на предприятии рабочей силы вместе с общим количеством затрат труда сохранилось, производство было полностью реорганизовано так, что на \textit{участке-1} стало трудиться 4 рабочих, на \textit{участке-2} и \textit{участке-3} по двое рабочих, количество производимых конечных изделий увеличилось с 20 до 40, экономический эффект составил 1280 человеко-часов.

Денежное выражение экономического эффекта по формуле:

$$E_v = v \times E_t = 5 \quad\mbox{тыс. руб.} \times 1280 \quad\mbox{человеко-часов} = 6400 \quad\mbox{тыс. руб.}$$

До внедрения технологии предприятие выпускало за месяц (20 рабочих дней) трудом 8-ми рабочих 20 конечных изделий, величина добавленной стоимости в общем месячном продукте 8-ми рабочих составляла 6400 тыс. руб. После внедрения технологии эти же 8 рабочих выпускают 40 конечных изделий, т.~е. в два раза больше, чем до внедрения технологии. При неизменной рыночной цене изделия величина добавленной стоимости в месячном продукте этих же 8-ми рабочих удвоится и составит 6400 тыс. руб. × 2 = 12800 тыс. руб. Таким образом, экономический эффект от повышения производительности труда после внедрения технологии составит добавочную долю добавленной стоимости в денежном выражении 12800 тыс. руб. - 6400 тыс. руб. = 6400 тыс. руб. (та же сумма, что и результат формулы).

Так же, как и в форме экономии затрат рабочего времени, в денежном выражении экономический эффект в сценарии внедрения технологии расширением производства с сохранением затрат труда (6400 тыс. руб.) превысил в два раза эффект внедрения технологии сокращением затрат рабочей силы с сохранением выпуска продукта (3200 тыс. руб.). При сохранении затрат труда рабочие получают необходимые жизненные средства, как и до внедрения технологии, не уходят за пределы предприятия, не получают в распоряжение дополнительное свободное время. Весь экономический эффект попадает в распоряжение владельцев предприятия полностью — как замещающий продукт, создаваемый меньшими силами, так и продукт, созданный вновь примененным высвободившимся трудом.

\section*{\raggedright{Реализация сокращения затрат рабочего времени снижением цены продукта}}

При расширении производства предприятию необходимо наладить сбыт добавочной массы продукта. На растущем рынке оно вполне сможет продать весь новый увеличенный объем продукта по текущей средней рыночной цене. В том случае, если спрос на рынке ограничен его текущим размером, для того, чтобы продать добавочную массу продукта, производитель должен привлечь новых потребителей к себе от конкурентов. При плюс-минус равных потребительских свойствах конкурирующих продуктов, если оставить в стороне такие способы конкурентной борьбы, как поглощения, лоббирование, заградительные пошлины и т.~п., один из действенных способов конкурентной борьбы — предложение своего продукта по цене ниже, чем у конкурентов. Цена продукта на конкурентном немонополизированном рынке в значительной степени обусловлена или, по крайней мере, ограничена снизу текущими средними по отрасли издержками на производство продукта и некоторой устоявшейся в отрасли нормой прибыли. 

В том случае, если производители имеют доступ к одному и тому же оборудованию, одним и тем же сырым материалам, применяют примерно одинаковые технологии и работают с примерно одинаковой производительностью, ни у кого из них не будет действенной возможности понизить цену, по крайней мере, если оставить в стороне такие механизмы, как демпинг. В том случае, если один из производителей внедрит новую технологию, заметно повышающую его индивидуальную производительность труда, в его распоряжении появится добавочный ресурс — экономический эффект, одним из путей реализации которого может быть экономически обоснованное снижение цены продукта с сохранением или увеличением текущей прибыльности.

Вернемся к сценарию внедрения технологии с сохранением затрат труда при увеличении объема производимого продукта. После внедрения технологии и повышения производительности труда производство было организовано таким образом, что всего в общей сложности 8 рабочих, задействованных на всех участках производства, производят в рассматриваемый период (месяц — 20 рабочих дней) 40 конечных изделий, общий экономический эффект составил 1280 человеко-часов, в денежном выражении 6400 тыс. руб.

При неизменных рыночных ценах, обусловленных средней по отрасли производительностью труда, равной старой производительности предприятия, количество добавленной стоимости в месячном продукте 40 изделий в денежном выражении составляет 12800 тыс. руб. Количество добавленной стоимости в цене одного изделия, таким образом, составляет 12800 тыс. руб. / 40 = 320 тыс. руб., из них доля экономического эффекта, приходящаяся на одно изделие, составляет 6400 тыс. руб. / 40 = 160 тыс. руб.

Предприятие имеет возможность понизить рыночную цену одного изделия на любую величину в диапазоне размера доли приходящегося на него экономического эффекта, сохранив при этом уровень прибыльности не ниже, чем был до внедрения технологии. Если предприятие понизит цену одного изделия на всю величину эффекта 160 тыс. руб., оно получит существенное конкурентное преимущество на рынке и все еще будет получать старую прибыль такую же, как и все прочие игроки на рынке, однако не получит никакой добавочной прибыли, из которой оно могло бы, например, компенсировать первоначальные издержки на внедрение технологии и реорганизацию производства. Если предприятие понизит цену одного изделия, к примеру, на половину эффекта 160 тыс. руб. / 2 = 80 тыс. руб., оно также получит ценовое преимущество перед конкурентами, вместе с этим в его распоряжении всё еще останется добавочная прибыль — оставшаяся доля экономического эффекта 80 тыс. × 40 = 3200 тыс. руб.

Доля экономического эффекта, реализованная как понижение цены, в таком случае уйдет за пределы предприятия в распоряжение потребителей, не будет отражена в учетных системах предприятия, по крайней мере, в финальной денежной форме. Тем не менее, для предприятия такой шаг будет экономически целесообразным, т.~к. благодаря ему оно сможет обеспечить сбыт добавочного продукта, полученного благодаря повышению производительности, сверх обычной массы продукта, которую оно сбывало ранее на уже занятой доле рынка.

Интересно отметить, что с понижением цены продукта уменьшится количество добавленной стоимости, извлекаемой предприятием с рынка в денежной форме в обмен на продукт, и, как следствие, понизится также денежная метрика производительности труда, по крайней мере, она не отразит полностью рост производительности, который произошел в действительности. В масштабах экономики, однако, доля экономического эффекта, реализованная понижением цены, отразится в сокращении расходов, которые несли потребители на покупку этого продукта. Если продукт предназначен для конечных потребителей, то это отразится в повышении их благосостояния. Если продукт предназначен для других предприятий, например, оборудование, сырые материалы, энергия и т.~п., то это позволит им, купив его по более низкой цене, получить собственный экономический эффект, который, в потенциале и в конечном итоге все равно будет реализован понижением цены уже на их собственные продукты, также предназначенные для конечных потребителей.

В условиях простой конкуренции для того, чтобы попытаться сохранить свои позиции на рынке, конкурирующие предприятия будут вынуждены снижать собственные цены. При прочих равных условиях единственный способ этого добиться — внедрить рассмотренную или аналогичную технологию, повысив собственную производительность труда. Пространство для «маневра» ценой внутри достигаемого таким образом экономического эффекта будет сжиматься до тех пор, пока он окончательно не уйдет за пределы всех предприятий сегмента рынка в распоряжение общества. Новая повышенная производительность превратится в производительность, среднюю по отрасли, — отправную точку для внедрения новых технологий и повторения цикла повышения производительности на новом витке.

\section*{\raggedright{Сокращение стоимости амортизационной составляющей оборудования, покупаемого на рынке}}

Для того, чтобы повысить производительность труда на производстве, можно оптимизировать рабочие процессы, организовать рабочее пространство (сократить простои и потери времени), разработать и внедрить новую более оптимальную программу для станка с ЧПУ или внедрить новое более производительное оборудование. Организационную оптимизацию процессов и внедрение нового программного обеспечения можно рассматривать как «чистые» технологии, которые оказывают влияние на производительность на старом уже имеющемся в наличии оборудовании — на первый взгляд может показаться, что их вклад в экономический эффект происходит исключительно в пределах добавленной стоимости уменьшением затрат труда рабочих. В случае внедрения нового оборудования, повышающего производительность рабочих, сразу очевидно, что кроме оценки эффекта в области производства добавленной стоимости, в расчет нужно также брать стоимость тела нового станка.

Стоимость нового рассматриваемого оборудования может быть меньше, больше или равна стоимости оборудования, выполняющего аналогичные функции, и которое уже есть в наличии. Производительность нового оборудования на целевых операциях может быть меньше, больше или равна производительности имеющегося в наличии оборудования. В том случае, если цена нового станка ниже, а производительность труда на нем выше, можно ожидать положительный эффект по результатам его внедрения. В том случае, если цена нового станка выше, а производительность ниже, можно ожидать отрицательный экономический эффект. Во всех остальных случаях требуется специальный расчет.

Для того, чтобы оценить долю стоимости применяемого на производстве оборудования в стоимости производимого с его применением продукта, обычно используют механизм амортизации. Простой вариант — перенести полную стоимость оборудования равномерно на весь произведенный с его помощью продукт за некоторый период, например, за 3 года или 5 лет, при нормальной целевой эксплуатации или на всё количество продукта, которое составит его оценочный ресурс.

Пусть $L_1$ — стоимость старого оборудования, $L_2$ — стоимость нового. Пусть продукт $M_1$ — это весь продукт, который можно произвести на старом станке за весь период полной амортизации при целевой эксплуатации, $M_2$ — продукт, который можно произвести за полный период амортизации на новом станке. Продукт $m_1$ можно произвести за рассматриваемый период времени — производственный цикл, — при нормальной целевой эксплуатации старого станка, продукт $m_2$ можно произвести за рассматриваемый производственный цикл при целевой эксплуатации нового станка. Продукт $M_1$ можно вывести как $M_1 = m_1 \times i_1$, где $i_1$ — количество рассматриваемых производственных циклов в периоде полной амортизации старого оборудования. Продукт $M_2$ можно вывести как $M_2 = m_2 \times i_2$, где $i_2$ — количество рассматриваемых производственных циклов в периоде полной амортизации нового оборудования.

Стоимость $L_1 / M_1$ — величина амортизационной составляющей в единице продукта, производимого на старом оборудовании. Стоимость $L_2 / M_2$ — величина амортизационной составляющей в единице продукта, производимого на новом оборудовании. Величина $L_2 / M_2 \times m_2$ — доля амортизационной составляющей в месячном объеме продукта, производимом на новом оборудовании. Величина $L_1 / M_1 \times m_2$ — доля амортизационной составляющей в таком объеме продукта, который производится за месяц на новом станке, но если бы его произвели на старом станке.

Экономический эффект сокращения амортизационной составляющей за рассматриваемый период времени в таком случае составит:

$$E = (L_1 / M_1) \times m_2 - (L_2 / M_2) \times m_2$$

Из формулы, в частности, видно, что даже в том случае, если стоимость нового оборудования и старого оборудования равны, увеличение производительности, сопровождаемое выпуском большей массы продукта в течение всего периода амортизации, приведет к тому, что доля амортизационной составляющей в единице продукта уменьшится, таким образом будет достигнут положительный экономический эффект.

Вернемся к примеру с предприятием. Пусть предприятие до внедрения нового оборудования за производственный цикл (один месяц — 20 рабочих дней) выпускало $m_1 = 20$ конечных изделий. Производственная линия разбита на 3 участка: \textit{участок-1} выпускает одну \textit{деталь-1} в день ($m_1^{(1)} = 20$), \textit{участок-2} выпускает одну \textit{деталь-2} в день ($m_1^{(2)} = 20$), на \textit{участке-3} рабочие собирают за день из одной \textit{детали-1} и одной \textit{детали-2} одно конечное изделие ($m_1 = m_1^{(3)} = 20$).

Пусть стоимость старого оборудования на \textit{участке-1} $L_1^{(1)} = 14400$ тыс. руб. (14 млн. 400 тыс. руб.). Стоимость оборудования на \textit{участке-2} $L_1^{(2)} = 7200$ тыс. руб., стоимость оборудования на \textit{участке-3} $L_1^{(3)}= 7200$ тыс. руб.

Пусть полный амортизационный период оборудования на всех участках составляет 36 месяцев (3 года). При выпуске $m_1^{(1)} = 20$ \textit{деталей-1} в месяц общий выпуск \textit{деталей-1} за весь амортизационный период оборудования на \textit{участке-1} составит $M_1^{(1)} = 20 \times 36 = 720$ штук. Аналогичным образом, общий выпуск \textit{деталей-2} $M_1^{(2)} = 720$, конечных изделий $M_1^{(3)} = 720$ за весь период амортизации.

Величина амортизационной составляющей в одной \textit{детали-1}, производимой на старом оборудовании, составит $L_1^{(1)} / M_1^{(1)}$ = 14400 тыс. руб. / 720 = 20 тыс. руб. Величина амортизационной составляющей в стоимости одной \textit{детали-2} составит $L_1^{(2)} / M_1^{(2)}$ = 7200 тыс. руб. / 720 = 10 тыс. руб. Величина амортизационной составляющей, присоединяемой к конечному изделию на \textit{участке-3}, составит $L_1^{(2)} / M_1^{(2)}$ = 7200 тыс. руб. / 720 = 10 тыс. руб. Общая амортизационная составляющая в одном конечном изделии до внедрения нового оборудования, таким образом, будет:

$$L_1^{(1)} / M_1^{(1)} + L_1^{(2)} / M_1^{(2)} + L_1^{(3)} / M_1^{(3)} = 20 + 10 + 10 = 40 \quad\mbox{тыс. руб.}$$

В общей стоимости продукта, производимого на старом оборудовании за рассматриваемый цикл производства, величина амортизационной составляющей составит 40 тыс. руб. × 20 = 800 тыс. руб.

Теперь предприятие решило внедрить новое более производительное оборудование на \textit{участке-1}, позволяющее сократить затраты труда при выпуске \textit{деталей-1}. Пусть стоимость нового оборудования для \textit{участка-1} оказалась выше, чем применявшегося до сих пор, и составила $L_2^{(1)}$ = 21600 тыс. руб. Технологии производства на \textit{участке-2} и \textit{участке-3} остались без изменений.

Пусть предприятие решило внедрить оборудование, сохранив прежний выпуск продукта и сократив затраты труда (соответствует рассмотренному ранее второму сценарию внедрения технологии). После внедрения оборудования выпуск \textit{деталей-1} на \textit{участке-1} $m_2^{(1)} = m_1^{(1)}$ = 20, выпуск \textit{деталей-2} на \textit{участке-2} $m_2^{(2)} = m_1^{(2)}$ = 20, выпуск конечных изделий на \textit{участке-3} $m_2^{(3)} = m_1^{(3)}$ = 20. Общий выпуск \textit{деталей-1} за весь амортизационный период нового оборудования на \textit{участке-1} при целевом выпуске $m_2^{(1)}$ = 20 составит $M_2^{(1)}$ = 20 × 36 = 720 штук, аналогичным образом на \textit{участке-2} $M_2^{(2)}$ = 720 штук и \textit{участке-3} $M_2^{(3)}$ = 720 штук. Оборудование на \textit{участке-2} и \textit{участке-3} осталось прежнее и по составу и по количеству, его стоимость осталась без изменений: $L_2^{(2)} = L_1^{(2)}$ = 7200 тыс. руб., $L_2^{(3)} = L_1^{(3)}$ = 7200 тыс. руб.

Величина амортизационной составляющей в одной \textit{детали-1}, производимой на новом оборудовании, составит $L_2^{(1)} / M_2^{(1)}$ = 21600 тыс. руб. / 720 = 30 тыс. руб. Амортизационная составляющая на \textit{участке-2} и \textit{участке-3} останется без изменений и составит $L_2^{(2)} / M_2^{(2)}$ = 7200 тыс. руб. / 720 = 10 тыс. руб., $L_2^{(3)} / M_2^{(3)}$ = 7200 тыс. руб. / 720 = 10 тыс. руб.

Общая амортизационная составляющая в одном конечном изделии после внедрения нового оборудования с сохранением старого объема производства, таким образом, будет:

$$L_2^{(1)} / M_2^{(1)} + L_2^{(2)} / M_2^{(2)} + L_2^{(3)} / M_2^{(3)} = 30 + 10 + 10 = 50 \quad\mbox{тыс. руб.}$$

В общей стоимости продукта, производимого на новом оборудовании за рассматриваемый цикл производства, величина амортизационной составляющей составит 50 тыс. руб. × 20 = 1000 тыс. руб.

Экономический эффект сокращения амортизационной составляющей за рассматриваемый период времени на \textit{участке-1} составит:

\begin{align}
E^{(1)} &= (L_1^{(1)} / M_1^{(1)}) \times m_2^{(1)} - (L_2^{(1)} / M_2^{(1)}) \times m_2^{(1)} =\nonumber\\
        &= 20 \quad\mbox{тыс. руб.} \times 20 - 30 \quad\mbox{тыс. руб.} \times 20 = -200 \quad\mbox{тыс. руб.}\nonumber
\end{align}

Экономические эффекты на \textit{участке-2} и на \textit{участке-3} будут ожидаемо равны нулю:

\begin{align}
E^{(2)} &= (L_1^{(2)} / M_1^{(2)}) \times m_2^{(2)} - (L_2^{(2)} / M_2^{(2)}) \times m_2^{(2)} =\nonumber\\
        &= 10 \quad\mbox{тыс. руб.} \times 20 - 10 \quad\mbox{тыс. руб.} \times 20 = 0 \quad\mbox{руб.}\nonumber
\end{align}

\begin{align}
E^{(3)} &= (L_1^{(3)} / M_1^{(3)}) \times m_2^{(3)} - (L_2^{(3)} / M_2^{(3)}) \times m_2^{(3)} =\nonumber\\
        &= 10 \quad\mbox{тыс. руб.} \times 20 - 10 \quad\mbox{тыс. руб.} \times 20 = 0 \quad\mbox{руб.}\nonumber
\end{align}

Общий экономический эффект сокращения амортизационной составляющей на всех участках:

$$E = E^{(1)} + E^{(2)} + E^{(3)} = -200 + 0 + 0 = -200 \quad\mbox{тыс. руб.}$$

Таким образом, величина амортизационной составляющей в единице конечного продукта при внедрении более дорогого оборудования с сохранением старого объема выпуска увеличилась. Экономический эффект сокращения амортизационной составляющей оказался отрицательным, сокращение обернулось ростом.

Однако, экономический эффект от сокращения затрат рабочего времени, посчитанный ранее, в этом сценарии в деньгах составил 3200 тыс. руб. за рассматриваемый цикл производства (один месяц) или 3200 тыс. руб. / 20 = 160 тыс. руб. на одно конечное изделие. Суммарный экономический эффект от сокращения затрат рабочего времени и сокращения (здесь — увеличения) амортизационной составляющей в общей массе конечного продукта за рассматриваемый цикл производства составил 3200 тыс. руб. - 200 тыс. руб. = 3000 тыс. руб. В итоге, общий экономический эффект внедрения нового оборудования оказался положительным, т.~к. выгода от повышения производительности труда превысила потери от увеличения амортизационной составляющей в стоимости продукта из-за более высокой цены станка.

Еще один сценарий внедрения оборудования — сохранить прежние затраты труда, увеличив выпуск продукта (соответствует рассмотренному ранее третьему сценарию внедрения технологии). После внедрения оборудования выпуск \textit{деталей-1} на \textit{участке-1} $m_2^{(1)}$ = 40, выпуск \textit{деталей-2} на \textit{участке-2} $m_2^{(2)}$ = 40, выпуск конечных изделий на \textit{участке-3} $m_2^{(3)}$ = 40.

Общий выпуск \textit{деталей-1} за весь амортизационный период нового оборудования на \textit{участке-1} при целевом выпуске $m_2^{(1)}$ = 40 составит $M_2^{(1)}$ = 40 × 36 = 1440 штук. На прежнем (не изменившемся) оборудовании за весь его амортизационный период будет выпущено \textit{деталей-2} на \textit{участке-2} $M_2^{(2)}$ = 40 × 36 = 1440 штук, конечных изделий на \textit{участке-3} $M_2^{(3)}$ = 40 × 36 = 1440 штук.

Хотя на 2-м и 3-м участках производительность труда и технологии производства не изменились, на них был удвоен выпуск продукта и были удвоены затраты труда. Для того, чтобы этого добиться, предприятие могло также удвоить производственные мощности на этих участках, например, закупив и введя в строй дополнительное оборудование тех же самых моделей. В таком случае стоимость оборудования на \textit{участке-2} составит $L_2^{(2)} = L_1^{(2)}$ × 2 = 14400 тыс. руб., на \textit{участке-3} составит $L_2^{(3)} = L_1^{(3)}$ × 2 = 14400 тыс. руб.

Величина амортизационной составляющей в одной \textit{детали-1}, производимой на новом оборудовании, составит $L_2^{(1)} / M_2^{(1)}$ = 21600 тыс. руб. / 1440 = 15 тыс. руб. Величина амортизационной составляющей в стоимости одной \textit{детали-2} составит $L_2^{(2)} / M_2^{(2)}$ = 14400 тыс. руб. / 1440 = 10 тыс. руб. Величина амортизационной составляющей, присоединяемой к конечному изделию на \textit{участке-3}, составит $L_2^{(3)} / M_2^{(3)}$ = 14400 тыс. руб. / 1440 = 10 тыс. руб.

Общая амортизационная составляющая в одном конечном изделии после внедрения нового оборудования с увеличением объема производства, таким образом, будет:

$$L_2^{(1)} / M_2^{(1)} + L_2^{(2)} / M_2^{(2)} + L_2^{(3)} / M_2^{(3)} = 15 + 10 + 10 = 35 \quad\mbox{тыс. руб.}$$

В общей стоимости продукта, производимого на новом оборудовании за рассматриваемый цикл производства, величина амортизационной составляющей составит 35 тыс. руб. × 20 = 700 тыс. руб.

Экономический эффект сокращения амортизационной составляющей за рассматриваемый период времени на \textit{участке-1} составит:

\begin{align}
E^{(1)} &= (L_1^{(1)} / M_1^{(1)}) \times m_2^{(1)} - (L_2^{(1)} / M_2^{(1)}) \times m_2^{(1)} =\nonumber\\
        &= 20 \quad\mbox{тыс. руб.} \times 40 - 15 \quad\mbox{тыс. руб.} \times 40 = 200 \quad\mbox{тыс. руб.}\nonumber
\end{align}

Экономические эффекты на \textit{участке-2} и на \textit{участке-3} будут снова нулевыми:

\begin{align}
E^{(2)} &= (L_1^{(2)} / M_1^{(2)}) \times m_2^{(2)} - (L_2^{(2)} / M_2^{(2)}) \times m_2^{(2)} =\nonumber\\
        &= 10 \quad\mbox{тыс. руб.} \times 40 - 10 \quad\mbox{тыс. руб.} \times 40 = 0 \quad\mbox{тыс. руб.}\nonumber
\end{align}

\begin{align}
E^{(3)} &= (L_1^{(3)} / M_1^{(3)}) \times m_2^{(3)} - (L_2^{(3)} / M_2^{(3)}) \times m_2^{(3)} =\nonumber\\
        &= 10 \quad\mbox{тыс. руб.} \times 20 - 10 \quad\mbox{тыс. руб.} \times 20 = 0 \quad\mbox{тыс. руб.}\nonumber
\end{align}

Общий экономический эффект сокращения амортизационной составляющей на всех участках:

$$E_{\mbox{сум}} = E^{(1)} + E^{(2)} + E^{(3)} = 200 + 0 + 0 = 200 \quad\mbox{тыс. руб.}$$

Эффект для периода полной амортизации:

$$E_{\mbox{полн}} = 200 \quad\mbox{тыс.} \times 36 = 7200 \quad\mbox{тыс. руб.}$$

Таким образом, хотя стоимость нового оборудования оказалась выше стоимости старого, в сценарии с увеличением выпуска продукта величина амортизационной составляющей в единице конечного продукта уменьшилась. Это вполне очевидно, т.~к. при большем выпуске продукта одна и та же стоимость станка будет распределена на большее количество единиц продукта. В сценарии внедрения оборудования с сохранением объема выпуска продукта, хотя повышение производительности дало положительный экономический эффект, недостаточный выпуск не раскрыл полный потенциал нового оборудования. Максимальный объем продукта, который можно произвести на данным оборудовании, ограничен количеством часов в сутках и техническими характеристиками станка. Если станок применяется на производстве так, что он производит меньше продукта, чем может технически, в экономическом эффекте его применения раскрывается не весь его потенциал.

Таким образом, при применении одного и того же оборудования при разных стратегиях внедрения экономический эффект сокращения амортизационной доли может быть как положительным, так и отрицательным. В этом случае он оказался положительным и дополнил эффект повышения производительности труда.

Хотя производство на \textit{участке-2} и \textit{участке-3} тоже расширилось в два раза, с увеличением производственных мощностей цена оборудования тоже удвоилась, поэтому экономический эффект сокращения амортизационной составляющей на этих участках ожидаемо оказался равен нулю.

Рассмотрим еще один вариант реализации рассмотренного сценария с сохранением затрат труда и увеличением выпуска продукта. На \textit{участке-2} и \textit{участке-3} производительность труда и технологии производства не изменились, при этом выпуск продукта был увеличен в два раза и количество рабочих на каждом из участков увеличилось с 1 до 2-х. Каждый из рабочих работает восьмичасовой рабочий день, поэтому для того, чтобы направить на участок еще одного рабочего, не обязательно докупать еще один комплект оборудования. В том случае, если технические характеристики имеющегося в наличии оборудования позволяют, можно удвоить производство с удвоением затрат труда, организовав работу в две смены. Стоимость оборудования на \textit{участке-2} в таком случае останется прежней и составит $L_2^{(2)} = L_1^{(2)}$ = 7200 тыс. руб., на \textit{участке-3} останется прежней и составит $L_2^{(3)} = L_1^{(3)}$ = 7200 тыс. руб. Выпуск продукта за весь амортизационный период: \textit{деталей-2} на \textit{участке-2} $M_2^{(2)}$ = 1440 штук, конечных изделий на \textit{участке-3} $M_2^{(3)}$ = 1440 штук.

Величина амортизационной составляющей в одной \textit{детали-1}, производимой на новом оборудовании, останется $L_2^{(1)} / M_2^{(1)}$ = 15 тыс. руб. Величина амортизационной составляющей в стоимости одной \textit{детали-2} составит $L_2^{(2)} / M_2^{(2)}$ = 7200 тыс. руб. / 1440 = 5 тыс. руб. Величина амортизационной составляющей, присоединяемой к конечному изделию на \textit{участке-3}, составит $L_2^{(3)} / M_2^{(3)}$ = 7200 тыс. руб. / 1440 = 5 тыс. руб.

Общая амортизационная составляющая в одном конечном изделии будет:

$$L_2^{(1)} / M_2^{(1)} + L_2^{(2)} / M_2^{(2)} + L_2^{(3)} / M_2^{(3)} = 15 + 5 + 5 = 25 \quad\mbox{тыс. руб.}$$

В общей стоимости продукта, производимого на новом оборудовании за рассматриваемый цикл производства, величина амортизационной составляющей составит 25 тыс. руб. × 20 = 500 тыс. руб.

Экономический эффект на \textit{участке-1} останется таким же, как в предыдущем варианте:

\begin{align}
E^{(1)} &= (L_1^{(1)} / M_1^{(1)}) \times m_2^{(1)} - (L_2^{(1)} / M_2^{(1)}) \times m_2^{(1)} =\nonumber\\
        &= 20 \quad\mbox{тыс. руб.} \times 40 - 15 \quad\mbox{тыс. руб.} \times 40 = 200 \quad\mbox{тыс. руб.}\nonumber
\end{align}

Экономические эффекты сокращения амортизационной составляющей на \textit{участке-2} и на \textit{участке-3} будут:

\begin{align}
E^{(2)} &= (L_1^{(2)} / M_1^{(2)}) \times m_2^{(2)} - (L_2^{(2)} / M_2^{(2)}) \times m_2^{(2)} =\nonumber\\
        &= 10 \quad\mbox{тыс. руб.} \times 40 - 5 \quad\mbox{тыс. руб.} \times 40 = 200 \quad\mbox{тыс. руб.}\nonumber
\end{align}

\begin{align}
E^{(3)} &= (L_1^{(3)} / M_1^{(3)}) \times m_2^{(3)} - (L_2^{(3)} / M_2^{(3)}) \times m_2^{(3)} =\nonumber\\
        &= 10 \quad\mbox{тыс. руб.} \times 20 - 5 \quad\mbox{тыс. руб.} \times 20 = 200 \quad\mbox{тыс. руб.}\nonumber
\end{align}

Общий экономический эффект сокращения амортизационной составляющей на всех участках:

$$E_{\mbox{сум}} = E^{(1)} + E^{(2)} + E^{(3)} = 200 + 200 + 200 = 600 \quad\mbox{тыс. руб.}$$

Эффект для периода полной амортизации:

$$E_{\mbox{полн}} = 600 \quad\mbox{тыс.} \times 36 = 21600 \quad\mbox{тыс. руб.}$$

Таким образом, хотя на \textit{участке-2} и \textit{участке-3} новые технологии не были внедрены и производительность труда не была повышена, доля амортизационной составляющей оборудования этих участков в конечном продукте была сокращена. Простое увеличение загрузки имеющегося в наличии оборудования с пропорциональным увеличением затрат труда дало заметный экономический эффект сокращения амортизационной составляющей при нулевом эффекте повышения производительности. Причина в том, что одна и та же стоимость оборудования при увеличении нагрузки распределилась на большую массу продукта. Из этого также следует, что после внедрения «чистых» технологий повышения производительности таких, как улучшенные версии программного обеспечения для станков с ЧПУ, на имеющемся в наличии парке оборудования кроме прямого эффекта повышения производительности следует также рассмотреть эффект сокращения амортизационной составляющей оборудования — при некоторых сценариях внедрения технологии он может оказаться не нулевым. В рассмотренном случае такой эффект получилось достигнуть простой реорганизацией рабочего процесса с имеющимся в наличии оборудованием и рабочей силой без внедрения новых технологий.

\section*{\raggedright{Экономия энергии, сырых материалов и прочих компонент, покупаемых за пределами предприятия}}

Внедрение новой технологии в дополнение к имеющемуся в наличии аппаратному обеспечению или в составе нового оборудования может привести к тому, что потребление энергии, сырых материалов, компонент и т.~п. может уменьшиться в пересчете на единицу конечного продукта, при этом потребительные свойства продукта окажутся прежними. Это приведет к тому, что индивидуальная стоимость продукта на предприятии уменьшится — даже в том случае, если изменений производительности труда на участке не произошло. Выгода от сокращения затрат на сырые материалы в составе производимого продукта создаст отдельный экономический эффект.

Расход сырых материалов, как и их экономия, могут быть выражены в натуральной форме, форме денег или затрат рабочего времени. Натуральная форма (штуки, вес, объем, площадь, длина, затраты энергии и т.~п.) является первичной в точке применения материала для производства из него продукта. В том случае, если материал производится внутри предприятия на одном из предшествующих участках производства, предприятие будет обладать информацией о затратах рабочего времени для его производства и сможет вести учет расхода и экономии в этой форме. Если предприятие закупает материал на рынке, в его распоряжении будет оценка расхода и экономии в деньгах.

Пусть $S_1^{(i, j)}$ — расход $j$-го материала или энергии, отраженный в системе учета ресурсов предприятия ERP, на $i$-м участке производства до внедрения новой технологии. После внедрения новой технологии система учета ресурсов предприятия ERP отразит расход $j$-го материала на $i$-м участке величиной $S_2^{(i, j)}$. Выпуск продукта на $i$-м участке до внедрения технологии — $m_1^{(i)}$, выпуск продукта на этом участке после внедрения технологии — $m_2^{(i)}$.

Расход $j$-го материала в пересчете на единицу продукта на $i$-м участке производства в рассматриваемый цикл производства до внедрения технологии составит $S_1^{(i, j)} / m_1^{(i)}$. Величина $(S_1^{(i, j)} / m_1^{(i)}) \times m_2^{(i)}$ — расход $j$-го материала на $i$-м участке производства для такого объема продукта, который был бы произведен на этом участке за рассматриваемый цикл производства после внедрения технологии, но если бы он был произведен без применения новой технологии. Экономический эффект экономии $j$-го материала или энергии на $i$-м участке производства, таким образом, составит:

$$E^{(i, j)} = (S_1^{(i, j)} / m_1^{(i)}) \times m_2^{(i)} - S_2^{(i, j)}$$

Из формулы, в частности, видно, что в том случае, если после внедрения технологии объем производства не изменился, экономический эффект экономии материала составит простую разницу между расходованием материала на один и тот же объем продукта. В том случае, если объем производства изменился и экономия материала, действительно, произошла, чем больше будет выпуск продукта после внедрения технологии, тем больше будет величина общего экономического эффекта.

Для каждого отдельного $j$-го материала экономический эффект можно считать в натуральной форме. Расход материалов и энергии, закупленных за пределами предприятия за деньги, можно перевести в денежную форму простым пересчетом натурального количества в соответствующую величину закупочной цены. В таком случае получится посчитать суммарный экономический эффект экономии сырых материалов и энергии, закупаемых за пределами предприятия, потребляемых на $i$-м участке производства:

\begin{align}
E^{(i)} = \textstyle\sum_{j} E^{(i, j)}, j \quad\mbox{—} \quad&\mbox{индексы наименований сырых материалов, энергии}\nonumber\\
                                                              &\mbox{и т.~п., закупаемых за пределами предприятия,}\nonumber\\
                                                              &\mbox{потребляемых на $i$-м участке производства}\nonumber
\end{align}

И суммарный экономический эффект экономии сырых материалов и энергии, закупаемых за пределами предприятия, потребляемых на всей производственной линии, в конечном продукте:

$$E_{\mbox{сум}} = E^{(1)} + … + E^{(n)}, n \quad\mbox{— количество участков производства}$$

Пусть предприятие до внедрения технологии за производственный цикл (один месяц — 20 рабочих дней) выпускало $m_1$ = 20 конечных изделий. Производственная линия разбита на 3 участка: \textit{участок-1} выпускает одну \textit{деталь-1} в день ($m_1^{(1)} = 20$), \textit{участок-2} выпускает одну \textit{деталь-2} в день ($m_1^{(2)} = 20$), на \textit{участке-3} рабочие собирают за день из одной \textit{детали-1} и одной \textit{детали-2} одно конечное изделие ($m_1 = m_1^{(3)} = 20$). На \textit{участке-1} при изготовлении \textit{детали-1} применяется электрическая энергия и некоторый сырой материал, к примеру, металл. На \textit{участке-2} при изготовлении \textit{детали-2} применяется электрическая энергия и некоторый сырой материал, к примеру, пластик. На \textit{участке-3} при изготовлении конечного изделия в качестве сырых материалов выступают \textit{детали-1} и \textit{детали-2}, произведенные на предыдущих участках, и электричество.

Пусть за рассматриваемый производственный цикл расход электричества на \textit{участке-1} в денежной форме составляет $S_1^{(1, 1)}$ = 100 тыс. руб., расход сырого материала в натуральной форме $S_{1\mbox{(нат)}}^{(1, 2)}$ = 900 единиц (веса, объема, площади и т.~п.), в денежной форме этот объем представлен закупочной ценой $S_1^{(1, 2)}$ = 600 тыс. руб. На \textit{участке-2} расход электричества составляет $S_1^{(2, 1)}$ = 100 тыс. руб., расход сырого материала $S_1^{(2, 2)}$ = 400 тыс. руб.

На \textit{участке-3} одно конечное изделие рабочие собирают из \textit{детали-1} и \textit{детали-2}, однако они не закупаются на рынке, а производятся внутри предприятия. Поэтому здесь в списке сырых материалов на \textit{участке-3} остается только потребление электричества $S_1^{(3, 1)}$ = 100 тыс. руб.

Пусть предприятие внедрило технологию, которая экономит расход сырого материала на \textit{участке-1} в два раза при изготовлении прежнего объема продукта, расход электричества при этом остаётся прежним. После внедрение технологии объем производства на всех участках остался прежним: $m_2^{(1)} = m_1{(1)} = 20$, $m_2^{(2)} = m_1^{(2)}$ = 20, $m_2^{(3)} = m_1^{(3)}$ = 20. В таком случае $S_2^{(1, 1)} = S_1^{(1, 1)}$ = 100 тыс. руб. Новый расход сырого материала на \textit{участке-1} в натуральной форме $S_{2\mbox{(нат)}}^{(1, 2)} = S_{1\mbox{(нат)}}^{(1, 2)} / 2$ = 450 единиц, расход в деньгах $S_2^{(1, 2)} = S_1^{(1, 2)} / 2$ = 300 тыс. руб. На \textit{участке-2} и \textit{участке-3} расход материалов без изменений: $S_2^{(2, 1)} = S_1^{(2, 1)}$ = 100 тыс. руб., $S_2^{(2, 2)} = S_1^{(2, 2)}$ = 400 тыс. руб., $S_2^{(3, 1)} = S_1^{(3, 1)}$ = 100 тыс. руб.

Экономия \textit{материала-2} на \textit{участке-1} в натуральной форме составит:

\begin{align}
E_{\mbox{(нат)}}^{(1, 2)} &= (S_{1\mbox{(нат)}}^{(1, 2)} / m_1^{(1)}) \times m_2^{(1)} - S_{2\mbox{(нат)}}^{(1, 2)} =\nonumber\\
                          &= (900 \quad\mbox{единиц} / 20) \times 20 - 450 \quad\mbox{единиц} = 450 \quad\mbox{единиц}\nonumber
\end{align}

Экономия \textit{материала-2} на \textit{участке-1} в деньгах составит:

\begin{align}
E^{(1, 2)} &= (S_1^{(1, 2)} / m_1^{(1)}) \times m_2^{(1)} - S_2^{(1, 2)} =\nonumber\\
           &= (600 \quad\mbox{тыс. руб.} / 20) \times 20 - 300 \quad\mbox{тыс. руб.} = 300 \quad\mbox{тыс. руб.}\nonumber
\end{align}

Экономия электричества (\textit{материала-1}) на \textit{участке-1}, очевидно, будет равна нулю:

\begin{align}
E^{(1, 1)} &= S_1^{(1, 1)} / m_1^{(1)} \times m_2^{(1)} - S_2^{(1, 1)} =\nonumber\\
        &= (100 \quad\mbox{тыс. руб.} / 20) \times 20 - 100 \quad\mbox{тыс. руб.} = 0 \quad\mbox{руб.}\nonumber
\end{align}

Аналогичным образом нулю в натуральной форме и деньгах будет равна экономия всех прочих материалов на всех остальных участках: $E^{(2, 1)}$ = 0 руб., $E^{(2, 2)}$ = 0 руб., $E^{(3, 1)}$ = 0 руб.

Т.~к. экономия материалов уже представлена в деньгах, эффекты для отдельных материалов можно суммировать, экономический эффект внедрения технологии на \textit{участке-1} составит:

$$E^{(1)} = E^{(1, 1)} + E^{(1, 2)} = 300 \quad\mbox{тыс. руб.} + 0 \quad\mbox{руб.} = 300 \quad\mbox{тыс. руб.}$$

На остальных участках общий экономический эффект равен нулю: $E^{(2)}$ = 0 руб., $E^{(3)}$ = 0 руб.

Общий суммарный экономический эффект внедрения технологии, экономящей сырой материал на \textit{участке-1}, составит:


$$E = E^{(1)} + E^{(2)} + E^{(3)} = 300 \quad\mbox{тыс. руб.} + 0 \quad\mbox{руб.} + 0 \quad\mbox{руб.} = 300 \quad\mbox{тыс. руб.}$$

В том случае, если внедрение технологии, сокращающей расход сырых материалов, сопровождалось увеличением выпуска продукта, как в случае простого расширения производства без повышения производительности, так и если эта технология входила в состав нового внедряемого более производительного станка, экономический эффект от сокращения потребления сырых материалов вырастет пропорционально увеличенному объему производства.

\section*{\raggedright{Экономия энергии, сырых материалов и прочих компонент, производимых внутри предприятия}}

В том случае, если энергия или сырой материал создается на рассматриваемом предприятии, к примеру, генератор вырабатывает электричество, или некоторый компонент создается на одном участке производства в качестве целевого изделия и потребляется в качестве составного элемента на следующем участке, их потребление тоже может быть сэкономлено внедрением некоторой технологии. Их объем, потребление и экономия будут выражены в натуральной форме, могут быть выражены в форме затрат рабочего времени и не будут выражены в деньгах. С точки зрения производственного процесса как такового, нет никакой разницы, произведен ли потребляемый материал внутри предприятия или куплен за его пределами за деньги. Однако, на производство материалов внутри предприятия затрачиваются не только другие материалы и оборудование, но и добавленный труд рабочих на соответствующих участках производства. Экономия таких материалов не может быть выражена сразу в простой экономии денежных средств, поэтому при расчете экономического эффекта их не следует включать в список материалов, покупаемых за его пределами, а следует рассмотреть отдельно.

Пусть предприятие до внедрения технологии выпускает $m_1 = 20$ конечных изделий за один производственный цикл (один месяц — 20 рабочих дней). Производственная линия разбита на 3 участка: \textit{участок-1} выпускает \textit{детали-1}, \textit{участок-2} выпускает \textit{детали-2}, на \textit{участке-3} происходит сборка конечного изделия из одной \textit{детали-1} и одной \textit{детали-2}. Процесс производства на \textit{участке-3} организован таким образом, что при сборке конечного изделия в среднем одна из 2-х \textit{деталей-2}, имеющихся в распоряжении, повреждается и отправляется в брак. Таким образом, хотя в составе конечного изделия присутствует всего одна \textit{деталь-1} и одна \textit{деталь-2}, при текущей организации производственного процесса при создании одного конечного изделия в среднем потребляется одна \textit{деталь-1} и две \textit{детали-2}.

Производство при этом организовано следующим образом: \textit{участок-1} выпускает в день (8 часов) 1 \textit{деталь-1} приложением труда 6-ти рабочих, \textit{участок-2} выпускает в день 2 \textit{детали-2} трудом 2-х рабочих, на \textit{участке-3} один рабочий собирает одно конечное изделие из одной \textit{детали-1} и двух \textit{деталей-2} (при этом в среднем одна из двух \textit{деталей-2} идет в брак). В месяц (20 рабочих дней) выпуск продукта и затраты труда будут: на \textit{участке-1} продукт $m_1^{(1)}$ = 20 шт, затраты труда $T_1^{(1)}$ = 8 часов × 6 человек × 20 дней = 960 человеко-часов, на \textit{участке-2} продукт $m_1^{(2)}$ = 40 шт, затраты труда $T_1^{(2)}$ = 8 часов × 2 человека × 20 дней = 320 человеко-часов, на \textit{участке-3} продукт $m_1 = m_1^{(3)}$ = 20 шт, затраты труда $T_1^{(3)}$ = 8 часов × 1 человек × 20 дней = 160 человеко-часов.

На \textit{участке-1} и \textit{участке-2} потребляются электрическая энергия и сырые материалы, производимые за пределами предприятия и покупаемые за деньги. На \textit{участке-3} потребляется электрическая энергия, покупаемая за деньги, в качестве сырых материалов выступают производимые на предприятии \textit{детали-1} и \textit{детали-2}. Таким образом, система ERP отразит потребление сырых материалов, производимых на предприятии, за месяц в натуральной форме: $S_{1\mbox{(нат)}}^{(3, 2)}$ = 20 шт. — потребление \textit{деталей-1}, $S_{1\mbox{(нат)}}^{(3, 3)}$ = 40 шт. — потребление \textit{деталей-2}.

Добавленные затраты рабочего времени в месячном продукте (20 изделий) до внедрения технологии:

$$T_1 = T_1^{(1)} + T_1^{(2)} + T_1^{(3)} = 960 + 320 + 160 = 1440 \quad\mbox{человеко-часов}$$

Добавленные затраты рабочего времени в одном конечном изделии до внедрения технологии составят: $T_1 / m_1 = 1440 / 20 = 72$ человеко-часа.

Пусть теперь на производство внедрена новая технология, которая улучшает производственный процесс на \textit{участке-3} таким образом, что повреждение и отбраковка половины \textit{деталей-2} при сборке конечного изделия прекращается, при сборке одного конечного изделия потребляется в среднем одна \textit{деталь-1} и одна \textit{деталь-2}.

С точки зрения рабочих на \textit{участке-3} произошла экономия сырых материалов (\textit{деталь-2}) при производстве конечного изделия. Для них при этом нет никакой разницы, была ли произведена \textit{деталь-2} внутри предприятия или куплена за пределами предприятия за деньги. Однако, в действительности, \textit{детали-2} производятся внутри предприятия, а не покупаются на рынке за деньги. Поэтому сокращение их потребления на \textit{участке-3} не может быть просто так конвертировано напрямую в эффект сокращения денежных затрат.

В том случае, если предприятие внедрит технологию, но при этом оставит структуру производства без изменений, выпуск изделий и затраты труда на всех участках производства останутся без изменений: на \textit{участке-1} продукт $m_2^{(1)} = m_1^{(1)} = 20$ шт, затраты труда $T_2^{(1)} = T_1^{(1)} = 960$ человеко-часов, на \textit{участке-2} продукт $m_2^{(2)} = m_1^{(2)} = 40$ шт, затраты труда $T_2^{(2)} = T_1^{(2)} = 320$ человеко-часов, на \textit{участке-3} продукт $m_2 = m_2^{(3)} = m_1^{(3)} = 20$ шт, затраты труда $T_2^{(3)} = T_1^{(3)} = 160$ человеко-часов.

После внедрения технологии система ERP отразит потребление сырых материалов, производимых на предприятии, за месяц в натуральной форме (всё потребление на \textit{участке-3}): $S_2^{(3, 2)} = 20$ шт. — потребление \textit{деталей-1}, $S_{2\mbox{(нат)}}^{(3, 3)} = 20$ шт. — потребление \textit{деталей-2}.

Экономия потребления \textit{деталей-1} на \textit{участке-3} в натуральной форме составит:

\begin{align}
E_{\mbox{(нат)}}^{(3, 2)} &= (S_{1\mbox{(нат)}}^{(3, 2)} / m_1^{(3)}) \times m_2^{(3)} - S_{2\mbox{(нат)}}^{(3, 2)} =\nonumber\\
                            &= (20 \quad\mbox{шт.} / 20) \times 20 - 20 \quad\mbox{шт.} = 0 \quad\mbox{шт.}\nonumber
\end{align}

Экономия потребления \textit{деталей-2} на \textit{участке-3} в натуральной форме составит:

\begin{align}
E_{\mbox{(нат)}}^{(3, 3)} &= S_{1\mbox{(нат)}}^{(3, 3)} / m_1^{(3)} \times m_2^{(3)} - S_{2\mbox{(нат)}}^{(3, 3)} =\nonumber\\
                            &= 40 \quad\mbox{шт.} / 20 \times 20 - 20 \quad\mbox{шт.} = 20 \quad\mbox{шт.}\nonumber
\end{align}

Таким образом, внедрение технологии на \textit{участке-3} позволило сэкономить 20 \textit{деталей-2} в месяц.

В затратах рабочего времени эта экономия составит $E^{(3, 2)}$ = 0 человеко-часов (экономия на \textit{деталях-1}), $E${(3, 2)} = 160 человеко-часов (экономия на \textit{деталях-2}). Суммарная экономия рабочего времени по отдельным материалам на всех участках:

$$E_{\mbox{сум}} = E^{(3, 2)} + E^{(3, 3)} = 0 + 160 = 160 \quad\mbox{человеко-часов}$$

В то же время добавленные затраты рабочего времени как в месячном продукте, так и в единице продукта, вполне очевидно, не изменились:

$$T_2 = T_2^{(1)} + T_2^{(2)} + T_2^{(3)} = 960 + 320 + 160 = 1440 \quad\mbox{человеко-часов}$$

В единице продукта: $T_2 / m_2 = 1440 / 20 = 72$ человеко-часа.

Общий экономический эффект экономии затрат труда от внедрения технологии оказался равен нулю.

$$E = T_1 - T_2 = 0 \quad\mbox{человеко-часов}$$

Таким образом, суммарный эффект экономии сырых материалов, производимых на предприятии, и общий эффект экономии затрат труда в конечном продукте не совпали. Экономия потребления сырых материалов при выбранной стратегии внедрения технологии в конечном продукте никаким образом себя не проявила.

Дело в том, что хотя потребление \textit{деталей-2} на \textit{участке-3} сократилось, их выпуск на \textit{участке-2} остался прежним. При сокращении потребления и сохранении прежнего выпуска \textit{деталей-2} появляется невостребованный излишек произведенных \textit{деталей-2} — перепроизводство. Т.~к. затраты труда на производство лишних \textit{деталей-2} происходят внутри предприятия, эти затраты присоединяются к итоговым затратам труда в конечном изделии, участки производства рассогласованы, итоговой экономии не произошло. Для того, чтобы экономия \textit{деталей-2} на \textit{участке-3} проявила себя сокращением затрат труда в конечном изделии, после внедрения технологии следует снова согласовать весь производственный процесс с учетом новой потребности производимых на разных участках производства компонент.

В зависимости от обстоятельств предприятие может решить расширить выпуск конечных изделий так, чтобы потребить весь объем оказавшихся невостребованными \textit{деталей-2} (нарастить выпуск \textit{деталей-1} на \textit{участке-1} и конечных изделий на \textit{участке-3}), или сократить выпуск \textit{деталей-2} до потребностей текущего выпуска конечных изделий на \textit{участке-3}.

Пусть предприятие решило сократить выпуск \textit{деталей-2} до новой необходимой потребности — одна \textit{деталь-2} в день, 20 в месяц. Сокращение производства сопровождается сокращением затрат применяемого на \textit{участке-2} рабочего времени — из двух рабочих один окажется сокращен, или двое рабочих переходят на сокращенный рабочий день. Новый выпуск и затраты рабочего времени в таком случае будут: на \textit{участке-1} продукт $m_2^{(1)} = m_1^{(1)}$ = 20 шт, затраты труда $T_2^{(1)} = T_1^{(1)}$ = 960 человеко-часов, на \textit{участке-2} продукт $m_2^{(2)}$ = 20 шт, затраты труда $T_2^{(2)}$ = 8 часов × 1 человек × 20 дней = 160 человеко-часов, на \textit{участке-3} продукт $m_2 = m_2^{(3)} = m_1^{(3)}$ = 20 шт, затраты труда $T_2^{(3)} = T_1^{(3)}$ = 160 человеко-часов.

Потребление сырых материалов, производимых на предприятии, за месяц в натуральной форме такое же, как при предыдущей стратегии внедрения технологии: $S_2^{(3, 2)}$ = 20 шт. (потребление \textit{деталей-1}), $S_{2\mbox{(нат)}}^{(3, 3)}$ = 20 шт. (потребление \textit{деталей-2}).

Экономия в натуральной форме как в предыдущей стратегии: $E_{\mbox{(нат)}}^{(3, 2)}$ = 0 шт., $E_{\mbox{(нат)}}^{(3, 3)}$ = 20 шт. Экономия в затратах рабочего времени: $E^{(3, 2)}$ = 0 человеко-часов (экономия на \textit{деталях-1}), $E^{(3, 2)}$ = 160 человеко-часов (экономия на \textit{деталях-2}). Суммарная экономия рабочего времени по отдельным материалам на всех участках:

$$E_{\mbox{сум}} = E^{(3, 2)} + E^{(3, 3)} = 0 + 160 = 160 \quad\mbox{человеко-часов}$$

Добавленные затраты рабочего времени в месячном продукте (20 изделий) после внедрения технологии экономии \textit{деталей-2} с сокращением их производства:

$$T_2 = T_2^{(1)} + T_2^{(2)} + T_2^{(3)} = 960 + 160 + 160 = 1280 \quad\mbox{человеко-часов}$$

Добавленные затраты рабочего времени в одном конечном изделии до внедрения технологии составят: $T_1 / m_1$ = 1280 / 20 = 64 человеко-часа.

Экономический эффект внедрения технологии, сопровождаемого реорганизацией производства, на одном цикле производства (один месяц):

$$E = T_1 - T_2 = 1440 - 1280 = 160 \quad\mbox{человеко-часов}$$

Как видно, в произведенной массе конечного продукта благодаря внедрению технологии и реорганизации производства произошло сокращение затрат труда. Благодаря внедрению технологии удалось достичь экономии деталей на \textit{участке-3}. Благодаря реорганизации производства общий конечный эффект экономии совпал с суммарным эффектом экономии материалов на отдельных участках — эффект экономии материалов отразился в общей экономии труда в конечном продукте.

Попробуем посчитать эффекты повышения производительности труда на отдельных участках производства:

На первом участке при производстве \textit{детали-1}:

\begin{align}
E^{(1)} &= (T_1^{(1)} / m_1^{(1)}) \times m_2^{(1)} - T_2^{(1)} =\nonumber\\
        &= (960 / 20) \times 20 - 960 = 0 \quad\mbox{человеко-часов}\nonumber
\end{align}

На втором участке при производстве \textit{детали-2}:

\begin{align}
E^{(2)} &= (T_1^{(2)} / m_1^{(2)}) \times m_2^{(2)} - T_2^{(2)} =\nonumber\\
        &= (320 / 40) \times 20 - 160 = 0 \quad\mbox{человеко-часов}\nonumber
\end{align}

На последнем участке при сборке конечного изделия:

\begin{align}
E^{(3)} &= (T_1^{(3)} / m_1^{(3)}) \times m_2^{(3)} - T_2^{(3)}\nonumber\\
        &= (160 / 20) \times 20 - 160 = 0 \quad\mbox{человеко-часов}\nonumber
\end{align}

Суммарный экономический эффект повышения производительности на всех участках:

$$E_{\mbox{сум}} = E^{(1)} + E^{(2)} + E^{(3)} = 0 + 0 + 0 = 0 \quad\mbox{человеко-часов}$$

Как видно, прямого повышения производительности труда не произошло ни на одном из участков — как на \textit{участке-2}, на котором было сокращено производство, так и на \textit{участке-3}, на котором сократилось потребление \textit{деталей-2}. Простое сокращение или расширение производства с пропорциональным сокращением или увеличением количества примененной рабочей силы не оказывает влияния на производительность труда, поэтому не создает экономический эффект. Сокращение потребления энергии, сырых материалов и прочих компонент при этом тоже не оказывает влияния на количество затрат труда, добавляемых к продукту рабочими рассматриваемого участка производства. Каждая отдельная \textit{деталь-1}, \textit{деталь-2} и конечное изделие получают прежнее количество добавленного труда на каждом из своих участков производства. Но итоговое количество добавленных затрат труда в конечном изделии оказывается меньше, т.~к. после внедрения технологии и реорганизации производства в его стоимость входит вместо двух всего одна \textit{деталь-2}. Кроме затрат рабочего времени в конечном продукте сократится потребление энергии, сырых материалов и оборудования, потреблением которых были произведены сэкономленные \textit{детали-2}.

Аналогичного эффекта можно было бы добиться упрощением конструкции конечного изделия с сохранением всех его потребительных качеств. К примеру, если бы при сборке конечного изделия никакие \textit{детали-2} не уходили в брак, но исходная конструкция включала бы две \textit{детали-2} в штатной комплектации, а после изменений в конструкции получилось бы использовать всего одну \textit{деталь-2} при идентичном наборе функций конечного изделия.

Таким образом, экономический эффект внедрения технологии, экономящей энергию и сырые материалы, производимые внутри предприятия, выражается в экономии общих итоговых добавленных затрат труда в конечном продукте. Кроме уменьшения доли добавленных затрат труда, в конечном продукте также сократится доля стоимости энергии, сырых материалов и амортизационной составляющей оборудования, приобретаемых за пределами предприятия и входящих в полную стоимость сэкономленных \textit{деталей-2}. Экономия достигается при условии, что выпуск на всех участках производства будет вновь согласован в соответствии с новыми потребностями производственного процесса.

Для того, чтобы реализовать экономический эффект экономии на сырых материалах, приходящих извне, достаточно сократить их закупку на очередном цикле производства. Для того, чтобы реализовать экономический эффект экономии сырых материалов, производимых на предприятии, необходимо перестроить производственную линию на новое потребление её отдельных частей. Стоит отметить, что по аналогичным причинам экономия энергии и материалов, закупаемых за пределами предприятия у внешних производителей, действительно реализованная внутри предприятия, не обязательно обернется экономией ресурсов в масштабах всей экономики. Для того, чтобы это произошло, их производитель-продавец должен сократить производство в соответствии с сокращением их потребления или обеспечить для сэкономленного объема ресурса дополнительный сбыт.

\section*{\raggedright{Сокращение непроизводственных издержек}}

Сотрудники предприятия, не принимающие непосредственного участия в производственном процессе, выполняют задачи, необходимые для нормальной работы и развития предприятия, — обеспечивают целевую производительность труда на предприятии через управление, планирование, учет. В системах учета, применяемых на предприятиях, сотрудников, выполняющих такие функции, могут относить к категориям ИТР (инженерно"=технические работы), служащие, руководство и т.~п.

Расходы на необходимые непроизводственные издержки включают зарплату сотрудников, выполняющих эти функции, расходы на офис, офисную технику, настольные компьютеры, ноутбуки гаджеты, электричество, расходные материалы такие, как бумага, канцелярия и т.~п. Предприятие компенсирует эти расходы из добавленной стоимости, создаваемой рабочими предприятия в сегменте производства.

Прямым эффектом внедрения технологий, повышающих производительность труда сотрудников сегмента необходимых непроизводственных издержек, является сокращение размера этих издержек и, таким образом, увеличение доли прибыли, которая остается в распоряжении владельцев предприятия. Вместе с тем, появление более совершенных инструментов в руках сотрудников, обеспечивающих необходимый уровень производительности труда на предприятии, может позволить им реализовать мероприятия и выстроить производственные процессы на более высоком уровне производительности труда. Таким образом, повышение производительности в непроизводственном сегменте в некоторых случаях может создать условия для повышения производительности на производстве, перехода предприятия на новый уровень целевой производительности, т.~е. к реализации вторичных производственных экономических эффектов.

К примеру, внедрение цифровой системы учета производства приведет к тому, что эту задачу будет решать меньшее количество человек, — это прямой эффект сокращения непроизводственных издержек. Кроме того, внедрение более совершенной системы учета позволит управляющему персоналу упорядочить производственные процессы, сократить простои оборудования, снизить потери и т.~п., т.~е. повысить производительность труда на производстве, что, в свою очередь, создаст дополнительные экономические эффекты в производственном сегменте.

Количество непроизводственных издержек будет зависеть от производственных процессов, выстроенных на предприятии, в частности, от количества сотрудников, от структуры и объема продукта, который оно выпускает. Таким образом, размер непроизводственных издержек на выбранном производственном цикле можно записать как некоторую функцию от затрат труда и объема производства:

$$C = C(T, m, o)$$

Где $T$ — затраты труда, $m$ — объем продукта на заданном периоде производства, $o$ — прочие факторы, например, внешние по отношению к предприятию обстоятельства.

Вполне очевидно, что с ростом объема производства и ростом количества занятых на производстве сотрудников будут расти также и непроизводственные издержки учета и управления задействованными ресурсами. Однако, в общем случае эта зависимость не обязательно будет линейной. К примеру, простое наращивание выпуска однородного продукта не обязательно приведет к заметной нагрузке на ведение типовых учетных операций. Тогда как расширение выпускаемой номенклатуры небольшим экспериментальным тиражом усложнит структуру учета и заметно увеличит нагрузку на учет и управление. Вывести универсальную конкретную форму зависимости непроизводственных издержек от структуры производства может быть затруднительно, если вообще возможно.

Вместе с тем, конкретную относительно достоверную форму зависимости можно попытаться вывести для конкретного предприятия, особенно, если производственные и операционные процессы в нем «устоялись» и повторяются от цикла к циклу. Количество необходимого учетного труда будет зависеть от набора и количества повторяющихся учетных операций, затраты времени на которые могут быть нормированы через механизмы снимков рабочего времени или определены статистически на основе накопленных исторических данных учета рабочего времени сотрудников непроизводственного сегмента. Затраты на отдельные нетиповые задачи, требующие творческого подхода или зависящие от внешних факторов за пределами предприятия, могут быть оценены экспертно и на основе исторических данных, в любом случае для них будут определены некоторые «разумные» рамки. Часть таких задач, если они повторяются от цикла к циклу из раза в раз, могут перейти в разряд нормированных по мере накопления предприятием опыта.

Пусть $C_1$ — непроизводственные издержки предприятия за рассматриваемый производственный цикл до внедрения некоторой технологии, повышающей производительность части сотрудников непроизводственного сегмента. Величина $C_2$ — непроизводственные издержки за такой же цикл после внедрения этой технологии.

Прямой экономический эффект внедрения этой технологии на рассматриваемом цикле производства в таком случае составит простое сокращение издержек:

$$E = C_1 - C_2$$

В том случае, если новый уровень издержек $C_2$ окажется меньше старого уровня издержек $C_1$, эффект $E$ будет положительным, внедрение технологии принесет предприятию выгоду. В том случае, если новый уровень издержек $C_2$ окажется выше старого уровня издержек $C_1$, расходы предприятия увеличатся.

В том случае, если внедрение новых инструментов управления и учета позволит сотрудникам непроизводственного сегмента внести улучшения в производственные процессы так, что будет повышена производительность в сегменте производства, эффекты реализации этих мероприятий следует вычислить дополнительно по методам расчета эффектов внедрения технологий, оказывающих влияние на производственный процесс.

В том случае, если технологии повышения производительности в непроизводственном сегменте не внедрялись, но были изменены структура, затраты труда и объем выпуска в сегменте производства, к примеру, после внедрения в него новых производственных технологий, количество непроизводственных издержек в таком случае все равно следует пересчитать, т.~к. их количество так или иначе зависит от текущей структуры производства. Реализация внедрения новой технологии на производстве через увеличение объема выпускаемого продукта может привести к росту непроизводственных издержек, т.~е. к отрицательному эффекту в непроизводственном сегменте. Однако в том случае, если положительный эффект роста производительности на производстве превысит отрицательный эффект в непроизводственном сегменте, внедрение такой технологии может оказаться все равно оправданным.

Пусть предприятие за рассматриваемый цикл (один месяц — 20 рабочих дней) производит 20 конечных изделий в общей сложности трудом 8-ми рабочих, распределенных по 3-м участкам производства. Пусть каждый рабочий присоединяет к продукту за месяц добавленной стоимости в денежном выражении 800 тыс. руб., при этом получает заработную плату за этот же период 200 тыс. руб. Таким образом, в общем объеме продукта, выпущенном за месяц (20 изделий), количество добавленной стоимости в деньгах составит 800 тыс. руб. × 8 = 6400 тыс. руб. (6 млн. 400 тыс. руб.). Из них фонд оплаты труда (ФОТ) рабочих на производстве составит 200 тыс. руб. × 8 = 1600 тыс. руб. Оставшаяся доля добавленной стоимости 6400 тыс. руб. - 1600 тыс. руб. = 4800 тыс. руб. — прибавочная стоимость, прибыль (маржинальная прибыль) предприятия. Часть её будет потрачена на обеспечение нормальной текущей деятельности предприятия — оплату непроизводственных издержек.

Пусть на предприятии также работают один бригадир, обеспечивающий надзор и контроль за целевой производительностью рабочих на производстве, два сотрудника, работающих с системой учета, один бухгалтер и один директор, выполняющий прочие функции планирования, принятия решений и контроля. Таким образом, в сегменте непроизводственных издержек предприятия занято всего 5 сотрудников. Пусть зарплата каждого из перечисленных сотрудников составляет всё те же 200 тыс. руб. за рассматриваемый цикл производства (месяц). ФОТ сегмента непроизводственных издержек предприятия, таким образом, составит $C_{1\mbox{(фот)}}$ = 200 тыс. руб. × 5 = 1000 тыс. руб. (1 млн. руб.). Прочие непроизводственные издержки пусть будут $C_{1\mbox{(пр.)}}$ = 300 тыс. руб. за этот же период. Все непроизводственные издержки до внедрения технологии:

\begin{align}
C_1 &= C_{1\mbox{(фот)}} + C_{1\mbox{(пр.)}} =\nonumber\\
    &= 1000 \quad\mbox{тыс. руб.} + 300 \quad\mbox{тыс. руб.} = 1300 \quad\mbox{тыс. руб.}\nonumber
\end{align}

Пусть предприятие внедрит новую технологию, к примеру, модуль цифровой системы учета и оборота документов, разработанный с учетом особенностей организации процессов на этом конкретном предприятии, который благодаря глубокой специализации и автоматизации наиболее частых типовых операций повысил производительность труда сотрудников, работающих с этой системой, в два раза. Объем операций, который до внедрения технологии выполняли двое сотрудников, работая полный рабочий день, после внедрения технологии сможет выполнить один сотрудник при полной занятости или два сотрудника при частичной занятости, работая на половине ставки.

Предприятие в таком случае сможет сократить одного из двух сотрудников, чья производительность была повышена, или перевести их обоих на половину ставки на неполный рабочий день. Так или иначе, ФОТ сотрудников сегмента непроизводственных издержек после внедрения технологии составит зарплату за 4 полные ставки $C_{2\mbox{(фот)}}$ = 200 тыс. руб. × 4 = 800 тыс. руб. Прочие непроизводственные издержки останутся прежними $C_{2\mbox{(пр.)}} = C_{1\mbox{(пр.)}}$ = 300 тыс. руб. Все непроизводственные издержки после внедрения технологии:

\begin{align}
C_2 &= C_{2\mbox{(фот)}} + C_{2\mbox{(пр.)}} =\nonumber\\
    &= 800 \quad\mbox{тыс. руб.} + 300 \quad\mbox{тыс. руб.} = 1100 \quad\mbox{тыс. руб.}\nonumber
\end{align}

Экономический эффект внедрения технологии, сокращающей непроизводственные издержки, таким образом, составит:

$$E = C_1 - C_2 = 1300 \quad\mbox{тыс. руб.} - 1100 \quad\mbox{тыс. руб.} = 200 \quad\mbox{тыс. руб.}$$

Вполне очевидно, в этом случае он составил величину зарплаты за полную ставку одного сотрудника, на которую уменьшили ФОТ сегмента непроизводственных издержек.

Структура производственной линии, производительность труда рабочих и объем выпуска конечного продукта при этом остались без изменений, таким образом осталась без изменений полная величина добавленной стоимости и величина прибыли, рассчитанная как добавленная стоимость за вычетом ФОТ рабочих (4800 тыс. руб.). До внедрения технологии доля прибыли после вычета ФОТ сотрудников сегмента непроизводственных издержек составляла $4800 \quad\mbox{тыс. руб.} - 1300 \quad\mbox{тыс. руб.}$ = 3500 тыс. руб. После внедрения технологии доля прибыли после вычета ФОТ сотрудников сегмента непроизводственных издержек составит $4800 \quad\mbox{тыс. руб.} - 1100 \quad\mbox{тыс. руб.}$ = 3700 тыс. руб. Добавочные 200 тыс. руб. — размер полученного экономического эффекта сокращения непроизводственных издержек, — перейдут в распоряжение владельцев предприятия и могут быть использованы ими по своему усмотрению, к примеру, реинвестированы в создание, закупку и внедрение новых технологий или получение добавочного дохода.

\section*{\raggedright{Оценка эффективности цифровой трансформации предприятия}}

В том случае, если одна технология проявляет себя сразу в нескольких точках производственной и организационной структуры, или если происходит одномоментное внедрение комплекса разнонаправленных технологий, общий суммарный экономический эффект, реализованный на предприятии на рассматриваемом периоде времени, сложится из отдельных элементов:

$$E = E_v + E_l + E_s + E_c$$

Здесь $E_v$ — экономия добавленного труда, включая экономию, достигнутую благодаря экономии энергии и сырых материалов, производимых внутри предприятия. $E_l$ — экономия амортизационной составляющей оборудования, покупаемого на рынке. $E_s$ — экономия энергии и сырых материалов, покупаемых на рынке. $E_c$ — экономия непроизводственных издержек.

Отдельные эффекты могут быть выражены в натуральной форме, форме затрат рабочего времени или денежной форме. В том случае, если эффекты представлены в натуральной форме, общий эффект может быть представлен как перечисление разнородных элементов. Для того, чтобы получить итоговый эффект в виде одного числа, все составляющие его эффекты должны быть представлены одновременно в форме затрат рабочего времени или одновременно в форме денег.

Эффекты, отражающие внутренние процессы предприятия, — экономия затрат труда на производстве и часть непроизводственных издержек, представляющая труд сотрудников, — имеют исходную форму затрат рабочего времени, но также могут быть конвертированы в деньги. Эффекты, связанные с расходованием внешних по отношению к предприятию ресурсов, — экономия амортизационной составляющей оборудования, экономия энергии и сырых материалов, покупаемых на рынке, — для предприятия представлены изначально в денежной форме. Хотя оборудование, энергия и прочие ресурсы имеют в своей основе затраты труда, приложенного их производителями, предприятие, покупающее их на рынке за деньги, не осведомлено о количестве этих затрат, в любом случае этот труд затрачивается за пределами предприятия-покупателя. В масштабах экономики можно получить оценку затрат труда для любого продукта, внутри предприятия это может быть затруднительно. Поэтому наиболее удобный вариант привести эффекты к единому виду — конвертировать их все в денежную форму.

Эффект прироста добавленной стоимости $E_v$ можно взять в первоначальной денежной форме с применением обобщенного дефлятора для отрасли. Но лучше взять как расчетный вариант, выведенный из предварительно вычисленного эффекта экономии затрат труда $E_v = v \times E_t$. В этом случае он более точно отразит повышение производительности, т.~к. на него не повлияет возможное целенаправленное или вынужденное снижение продажной цены на продукт на рынке, результат будет равноценен применению дефлятора, вычисленного для индивидуальной ценовой политики предприятия. Эффект экономии затрат труда $E_t$ будет включать как эффект прямого повышения производительности труда на отдельных участках производства, так и эффект экономии энергии и сырых материалов, производимых внутри предприятия.

Для того, чтобы реализовать экономические эффекты, предприятие несет издержки на покупку, разработку и внедрение технологии, в общей сложности эти расходы составят некоторую величину — цену внедрения $P$. В том случае, если общий эффект внедрения $E$ оказался положительным, для оценки целесообразности внедрения технологии на целевом периоде времени следует его сравнить с ценой внедрения $P$, вычислив разницу $(E - P)$.

Если эффект $E$ больше издержек $P$, $(E - P) > 0$, повышенная производительность окупит понесенные издержки уже на рассмотренном периоде времени, внедрение можно считать целесообразным. Если эффект $E$ меньше либо равен издержкам $P$, $(E - P) \le 0$, повышенная производительность выведет в ноль или не окупит понесенные расходы. Следует отметить, что на более длинном периоде времени эффект внедрение $Е$ будет расти, т.~е. накапливаться, исходный размер издержек $P$ при этом останется без изменений. Поэтому на достаточно длинном периоде времени внедрение технологии всегда окажется целесообразным даже в том случае, если на более коротком периоде это окажется не так, вопрос только в продолжительности этого периода.

Цена внедрения $P$ — безвозвратные расходы, которые компенсируются экономическим эффектом повышения производительности труда на предприятии, внедрившем технологию. Она будет включать такие расходы, как покупка лицензий на программное обеспечение, зарплата разработчиков программного обеспечения, инженеров, консультантов по организации производственного процесса (бережливое производства, система 5S и т.~п.), прочие издержки и расходы, связанные непосредственно с процедурой внедрения технологии и периодом реорганизации производства. В том случае, если внедрение технологии сопровождалось дополнительными инвестициями, к примеру, в расширение производства, закупку нового оборудования, привлечение дополнительной рабочей силы и т.~п., их величину не следует относить к цене внедрения, т.~к. инвестиции такого рода в итоге входят в стоимость продукта и составляют часть его цены или компенсируются из добавленной стоимости как необходимые непроизводственные издержки в рамках обновленного производственного процесса.

В том случае, если предприятие ведет подробный детализированный учет ресурсов в системах учета ресурсов ERP и затрат труда в системах учета рабочего времени, а также историю внедрения новых технологий, при этом маркирует историю затрат ресурсов, производства продуктов, затрат труда и т.~п., ассоциируя их с примененными в производственном процессе конкретными технологиями, оно сможет оценить экономические эффекты внедрения отдельных технологий или семейств технологий. Таким образом предприятие может оценить и сравнить эффективность внедрения и применения каждой из них.

\section*{\raggedright{Повышение производительности труда инженеров — сокращение расходов на инженерные разработки}}

Накапливаемый экономический эффект от повышения производительности труда раньше или позже окупит расходы на разработку, покупку и внедрение технологий, которые позволили достичь этого эффекта. Однако предприятие может быть ограничено в ресурсах — абсолютной величине свободных средств, необходимых для реализации проекта повышения производительности. Кроме того, предприятие может иметь свои требования к периоду окупаемости такого проекта, в течение которого накопленный экономический эффект полностью компенсирует первоначальные вложения — цену внедрения проекта.

Таким образом, снижение цены внедрения технологии, во-первых, сокращает период окупаемости, т.~к. при одном и том же уровне новой производительности и достигнутом эффекте, быстрее окупится такая технология, цена внедрения которой оказалась ниже. Во-вторых, более низкая цена внедрения расширяет круг предприятий, которые могут себе позволить внедрить такую технологию.

Чем шире растиражирована универсальная технология, тем ниже её относительная цена внедрения для каждого конечного потребителя. Цена внедрения индивидуализированных разработок для их потребителей будет максимально высока. Главный путь сокращения абсолютной величины цены внедрения технологии — повышение производительности труда разработчиков технологии и прочих специалистов, занятых в реализации проекта, — инженеров-разработчиков программного обеспечения, прочих инженеров-проектировщиков, прочих специалистов, консультантов и т.~п. Повышение производительности инженеров, кроме сокращения цены, может также сократить время самой разработки.

Пусть предприятие за производственный цикл в один месяц (20 рабочих дней) выпускает 20 конечных изделий трудом 8-ми рабочих, распределенных по участкам производства. Для того, чтобы узнать возможные пути повышения производительности, предприятие привлекает внешнюю инженерную компанию, специализирующуюся на разработке технологий автоматизации производства. В силу ряда причин владельцы предприятия поставили условие, что экономический эффект внедрения технологии должен компенсировать величину инвестиций — цены разработки и внедрения технологии, — не позднее, чем через 6 месяцев (6 циклов) после запуска реорганизованной линии.

Специалисты инженерной компании оценили производственные процессы предприятия и предложили вариант разработки и внедрения технологии, которая повысит производительность труда на одном из участков. Расчетный экономический эффект предложенного сценария внедрения составил 3200 тыс. руб. за цикл в один месяц, 3200 тыс. руб. × 6 = 19200 тыс. руб. (19 млн. 200 тыс. руб.) за максимально приемлемый цикл компенсации инвестиций (6 месяцев), благодаря которым будет разработана и внедрена технология, которая после внедрения позволит достигнуть расчетный экономический эффект.

Специалисты инженерной компании оценили фронт работ и решили, что для разработки и внедрения технологии до момента запуска обновленного производства в полном объеме, потребуется один год. При этом в проекте на условиях полной занятости будут задействованы специалисты: 4 инженера-разработчика программного обеспечения, 1 консультант, 1 руководитель — на стороне инженерной компании-исполнителя, еще 2 специалиста со стороны заказчика, всего 8 человек. При зарплате 200 тыс. руб. в месяц каждого из специалистов за год (12 месяцев) разработки и внедрения проекта общие затраты на оплату труда специалистов составят 200 тыс. руб. × 8 × 12 = 19200 тыс. руб. Еще 4000 тыс. руб. — покупка лицензий на стороннее программное обеспечение, операционные расходы, прочие расходы. Таким образом, расчетная сумма цены внедрения составила $P$ = 19200 тыс. руб. + 4000 тыс. руб. = 23200 тыс. руб. (23 млн. 200 тыс. руб.).

Хотя расчетный экономический эффект внедрения технологии за год после начала её эксплуатации составит $E_{\mbox{год}}$ = 3200 тыс. руб. × 12 = 38400 тыс. руб. и превысит цену внедрения $P$ = 19200 тыс. руб., по условиям, выставленным владельцами предприятия, срок компенсации инвестиций не должен превышать 6 месяцев. Расчетный потенциальный экономический эффект за 6 месяцев $E$ = 19200 тыс. руб. оказывается меньше цены внедрения $P$ = 23200 тыс. руб.

$$E - P = 19200 \quad\mbox{тыс. руб.} - 23200 \quad\mbox{тыс. руб.} = -4000 \quad\mbox{тыс. руб.} < 0$$

Таким образом, владельцы предприятия принимают решение не финансировать предложенный проект.

Пусть теперь инженерная компания решила внедрить уже в свои процессы некоторую технологию или комплекс технологий (программный фреймворк, среду разработки, набор библиотек, ИИ-ассистента и т.~п.), которая позволила, действительно, повысить производительность труда инженеров-разработчиков программного обеспечения, занятых в компании. При новой оценке объем работ, который до внедрения технологии требовал труда 4-х инженеров-разработчиков программного обеспечения в течение года, после внедрения технологии разработки может быть выполнен 2-мя инженерами-разработчиками тоже за год.

В той же самой исходной ситуации переговоров с предприятием с новыми вводными условиями специалисты инженерной компании теперь решили, что для реализации проекта повышения производительности на предприятии потребуется один год и будут задействованы специалисты: 2 инженера-разработчика программного обеспечения (вместо предыдущих четырёх), 1 консультант, 1 руководитель — на стороне инженерной компании-исполнителя, 2 специалиста со стороны заказчика, всего 6 специалистов. При зарплате 200 тыс. руб. в месяц каждого из специалистов за год разработки и внедрения проекта общие затраты на оплату труда специалистов составят 200 тыс. руб. × 6 × 12 = 14400 тыс. руб. Плюс 4000 тыс. руб. — покупка лицензий на стороннее программное обеспечение, операционные расходы, прочие расходы. Таким образом, расчетная сумма цены внедрения составила в этом случае $P$ = 14400 тыс. руб. + 4000 тыс. руб. = 18400 тыс. руб. (18 млн. 400 тыс. руб.).

Таким образом, при новом уровне производительности труда специалистов инженерной компании расчетный потенциальный экономический эффект за 6 месяцев $E$ = 19200 тыс. руб. оказывается больше цены внедрения $P$ = 14400 тыс. руб.

$$E - P = 19200 \quad\mbox{тыс. руб.} - 14400 \quad\mbox{тыс. руб.} = 4800 \quad\mbox{тыс. руб.} > 0$$

Владельцы предприятия принимают решение о запуске предложенного проекта, выделяют инвестиционные средства, которые рассчитывают компенсировать из добавочного дохода — реализации экономического эффекта на предприятии, — не позднее, чем через 6 месяцев после внедрения разработанной технологии, или не позднее 1 года и 6 месяцев после выделения инвестиций, запуска процесса разработки технологии и её последующего внедрения.

Из приведенного примера хорошо видно, что эффекты повышения производительности труда инженеров, в т.~ч. разработчиков программного обеспечения, получают действительную реализацию только после того, как их разработки завершены и внедрены в реальный производственный процесс. Повышение производительности инженеров, понижая цену разработки, уменьшает цену внедрения их продукта в производственный процесс, снижает порог целесообразности внедрения их продуктов для повышения производительности труда на предприятиях и реализации новых экономических эффектов. Внедрение разработок на предприятиях реализует экономические эффекты в реальном секторе, которые не были бы реализованы, если бы цена разработки и внедрения технологии была бы слишком высока. В этом заключается опосредованная роль инструментов повышения производительности труда инженеров-разработчиков в росте и повышении производительности всей экономики.

\section*{\raggedright{Особенность цифровых технологий при повышении производительности предприятий и организаций}}

Технологии, основанные на программируемых цифровых вычислениях, в настоящий момент являются доминирующим инструментом инженерных разработок. Есть основания утверждать, что любое или почти любое достигаемое повышение производительности труда сейчас — в той или иной степени заслуга цифровых технологий или происходит при их участии.

Сегмент непроизводственных издержек — учет, планирование, взаимодействие с контрагентами, сбыт, документооборот и т.~п., — за годы развития технологий цифровых вычислений в глубочайшей степени трансформировался в цифровую форму. Значительная часть задач в сегменте целиком или почти целиком осуществляется в цифровом пространстве. Технологические гиганты такие, как Microsoft, Google, Oracle, Amazon и т.~п., получают прибыль благодаря цифровой автоматизации тех или иных задач внутри сегмента непроизводственных издержек. Выбирая точки приложения новых технологий с целью сокращения издержек и повышения производительности труда, предприятия и организации зачастую выбирают как первоочередную цель трудовые функции сотрудников сегмента непроизводственных издержек \cite{minecoLabourProductivitySectorProgramsForMay2025}. Вероятно, не в последнюю очередь потому, что алгоритмы такого рода запускаются на унифицированных массовых аппаратных платформах — настольных компьютерах, ноутбуках, мобильных гаджетах, серверной аппаратуре, — и распространяются в форме «чистого» программного обеспечения или как сетевые веб-приложения. Централизация аппаратной инфраструктуры облачных сервисов автоматизации офиса и инструментов сбыта через размещение рекламы таких, как поисковые системы, соцсети, маркетплейсы, при этом, делают инструменты такого рода максимально доступными. Цена внедрения новой цифровой технологии в сегмент непроизводственных издержек, таким образом, в той части, где она зависит от технологий, стремится к минимуму и в абсолютном значении уже достаточно низка. Другая причина концентрации внимания на сегменте непроизводственных издержек может заключаться в том, что автоматизация таких задач оказывается наиболее простым вариантом реализации экономического эффекта — это простое сокращение расходов, которое может быть достигнуто уже на очередном цикле производства после внедрения технологии.

Распространение и внедрение инструментов автоматизации инженерных разработок, как и в случае с непроизводственными издержками, по критерию сложности доставки и развертывания является максимально простой процедурой. Для специализированных технологических компаний, ведущих деятельность по разработке интеллектуальных продуктов, в т.~ч. программного обеспечения, на постоянной цикличной основе, регулярно реинвестирующих ресурсы на исправление старых и выпуск новых версий продукта, реализация эффекта повышения производительности труда инженеров выглядит как простое сокращение расходов, как и в случае с сокращением непроизводственных издержек. Однако, действительный экономический эффект инженерной разработки создается в точке ее внедрения, т.~е. после того, как продукт уже разработан. Сокращение издержек на разработку проявляется не в получении собственниками добавочной доли прибыли, а в сокращении необходимого для разработки объема инвестиций.

В области производства технологии цифрового программирования наделяют оборудование высокой степенью автономности и универсальности, исключая необходимость участия человека в множестве промежуточных производственных операций, таким образом повышая производительность производственного труда. Как и в случае с универсальными цифровыми потребительскими устройствами — настольными компьютерами, ноутбуками, мобильными гаджетами, — загрузка цифрового программного обеспечения позволяет выполнять широкий спектр задач на одной и той же производственной аппаратной платформе, к примеру, на станке с ЧПУ, вносить изменения и улучшения в производственный процесс без полной смены парка оборудования.

Вместе с тем, темпы и масштабы внедрения новых цифровых технологий, например, нового программного обеспечения, в секторе производства могут быть не такими стремительными и широкими, как в случае с сектором непроизводственных издержек или с сектором инженерных разработок. Возможно, отчасти из-за того, что аппаратные производственные платформы унифицированы в гораздо меньшей степени, чем массовые потребительские вычислительные устройства. Отчасти из-за того, что производственные процессы внутри предприятий обычно закрыты, разнообразны и унифицированы в еще меньшей степени, чем применяемое производственное оборудование. Отчасти из-за того, что алгоритмы, выполняемые на производственном оборудовании, в гораздо большей степени зависят от контекста конкретного производственного процесса, например, применяемого в текущий момент сырого материала, в отличие от персональных гаджетов, контекст которых в значительной степени ограничен экраном, клавиатурой и некоторым количеством стандартизированных устройств ввода и вывода. Отчасти из-за того, что для внедрения нового оборудования в том случае, если его требуется закупить, требуется гораздо больше вложений капитала, чем для, к примеру, покупки офисной техники.

Кроме объективных технических факторов такого рода, внедрение технологий на производстве осложняется тем, что точечное повышение производительности на отдельном участке не обязательно превратится в экономический эффект сокращения стоимости конечного продукта, если темпы и объемы производства на отдельных участках производственной линии не будут вновь согласованы с учетом новой производительности. Сценарий внедрения технологии с сохранением выпуска и сокращением затрат труда в точке внедрения отчасти решает проблему необходимости пересогласования всей линии, но в таком случае не раскрывает полный потенциал реализации эффекта повышения производительности.

Более глубокая степень цифровизации производственных процессов открывает дорогу для реализации более выгодных стратегий повышения производительности. В том случае, если программируемое управляемое цифровыми алгоритмами оборудование во всех точках производственной линии объединено в управляемую сеть в рамках единого цифрового пространства, включающего инструменты мониторинга, учета времени и ресурсов, контроля и управления, применение системы такого рода позволит спланировать реорганизацию, перераспределить нагрузку, минимизировать время перенастройки оборудования при повышении производительности на отдельных участках производственной линии или при комплексной модернизации технологического процесса \cite{minecoLabourProductivityCoverage2024, minecoIntexVolumeExpansion2025, vedomostiRussianIndustryEfficiency2024}. Это, в итоге, уменьшит цену внедрения каждой новой технологии, таким образом понизит порог перехода на каждый новый уровень производительности.

В том случае, если система комплексной глубокой цифровизации превратит производство в единую цифровую среду производства и управления, таким образом опустит цену внедрения каждой новой технологии достаточно низко, это позволит запустить процесс регулярного внедрения новых технологий, постоянного повышения производительности в сегменте производства, реализации сопутствующих экономических эффектов таких, как снижение цен на продукты потребления, сокращение рабочего дня с пополнением фонда общественного времени досуга, получение добавочной прибыли.

% чтобы ссылки в тексте имели правильные номера, генератор нужно запустить два раза
% нижние подчеркивания '_' в ссылках заменяем на '\_', иначе они воспринимаются как спец-символы, если они не обернуты в \url{}
% если ссылка обернута в \url{}, нижние подчеркивания экранировать не надо
\begin{thebibliography}{2}
\bibitem{sapDeloitteSurvey_eg_2021} Цифровая зрелость российских компаний требует развития / Игорь Борисов // Экономика и жизнь, – 2021. [Электронный ресурс]. – URL: \url{https://www.eg-online.ru/article/436133/}
\bibitem{sapDeloitteSurvey_deloitte_2021} Исследование SAP и «Делойта»: как повысить уровень цифровой зрелости — от стратегии до внедрения // Deloitte, – 2021. [Электронный ресурс]. – URL: \url{https://web.archive.org/web/20210414213055/https://www2.deloitte.com/ru/ru/pages/consulting/articles/sap-deloitte-research.html}
\bibitem{bcgSurvey2021} Догнать цифровое будущее Как сократить цифровое отставание в металлургической и горнодобывающей промышленности // BCG, – 2021. [Электронный ресурс]. – URL: \url{https://web.archive.org/web/20220810023743/https://www.bcg.com/press/25february2021-catching-digital-future-to-close-the-digital-lag-metals-mining-industry}
\bibitem{networkReadinessIndex2024} NRI 2024 Edition – Press Release // Portulans Institute, – 2024. [Электронный ресурс]. – URL: \url{https://networkreadinessindex.org/nri-2024-edition-press-release/}
\bibitem{hseDigitalStats2024} Индикаторы цифровой экономики: 2024: статистический сборник / Абашкин~В.~Л., Абдрахманова~Г.~И., Вишневский~К.~О., Гохберг~Л.~М., Демидкина~О.~В., Демьянова~А.~В., Зинина~Т.~С., Зорина~О.~А., Ковалева~Г.~Г., Коцемир~М.~Н., Кузина~Л.~С., Кузнецова~И.~А., Лола~И.~С., Мартынов~Д.~М., Нестеренко~А.~В., Нечаева~Е.~Г., Озерова~О.~К., Остапкович~Г.~В., Покровский~С.~И., Попов~Е.~В., Ратай~Т.~В., Репина~А.~А., Рудник~П.~Б., Рыжикова~З.~А., Стрельцова~Е.~А., Фридлянова~С.~Ю., Филатов~М.~М., Фролов~М.~С., Шугаль~Н.~Б., Щербаков~Р.~А., Никитская~Е.~Д., под общей редакцией: Анисимов~Н.~Ю., Гохберг~Л.~М., Кузьминов~Я.~И., А.~В.~Чукарин, И.~В.~Васильев // НИУ ВШЭ, – 2024. [Электронный ресурс]. – URL: \url{https://publications.hse.ru/books/888949257}
\bibitem{rosstat_prikaz332_31072024_3Inform} Приказ №~332 от 31.07.2024 Об утверждении форм федерального статистического наблюдения для организации федерального статистического наблюдения за деятельностью в сфере образования, науки и инноваций, Приложение №~11 Форма №~3-информ // Федеральная служба государственной статистики, – 2024. [Электронный ресурс]. – URL: \url{https://rosstat.gov.ru/folder/12981}
\bibitem{rosstat_postanov158_25072002_3Inform} Постановление от 25 июля 2002 г. N~158 Об утверждении статистического инструментария для организации статистического наблюдения за деятельностью, осуществляемой в сфере услуг, транспорта и связи, жилищно-коммунального хозяйства, здравоохранения, правонарушений, науки и инноваций на 2003 год // Государственный комитет Российской Федерации по статистике, – 2002. [Электронный ресурс]. – URL: \url{https://normativ.kontur.ru/document?moduleId=1&documentId=8387}
\bibitem{rosstat_postanov39_04032003_3Inform} Постановление от 4 марта 2003 г. N~39 Об утверждении инструкции по заполнению формы федерального государственного статистического наблюдения N 3-информ "сведения об использовании информационных технологий и производстве связанной с ними продукции (работ, услуг)" // Государственный комитет Российской Федерации по статистике, – 2003. [Электронный ресурс]. – URL: \url{https://normativ.kontur.ru/document?moduleId=1&documentId=56661}
\bibitem{oecdLabourProductivityManual2001} Measuring Productivity - OECD Manual // OECD, – 2001. [Электронный ресурс]. – URL: \url{https://www.oecd.org/en/publications/measuring-productivity-oecd-manual_9789264194519-en.html}
\bibitem{iloLabourProductivityKILM2016} Key Indicators of the Labour Market, Ninth edition // ILO, – 2016. [Электронный ресурс]. – URL: \url{https://www.ilo.org/publications/key-indicators-labour-market-2015-kilm-full-report}
\bibitem{rosstat_prikaz274_28042018_LabourProductivityIndex} Приказ от 28 апреля 2018 г. №~274 Об утверждении Методики расчета показателя "Индекс производительности труда" // Федеральная служба государственной статистики, – 2018. [Электронный ресурс]. – URL: \url{https://www.garant.ru/products/ipo/prime/doc/71838328/}
\bibitem{rosstatLabourProductivityIndexData2025} Национальные счета, Индекс производительности труда // Федеральная служба государственной статистики, – 2025. [Электронный ресурс]. – URL: \url{https://rosstat.gov.ru/statistics/accounts}
\bibitem{mineco_prikaz817_23122024_LabourProductivityGrowth} Приказ от 23 декабря 2024 г. N~817 Об утверждении методики расчета показателя "Рост производительности труда в экономике Российской Федерации по отношению к 2023 году" // Министерство экономического развития Российской Федерации, – 2024. [Электронный ресурс]. – URL: \url{https://www.garant.ru/products/ipo/prime/doc/411124549/}
\bibitem{rosstat_prikaz76_16022018_LabourProductivityDynamics} Приказ от 16 февраля 2018 г. N~76 Об утверждении Методики расчета показателя "Динамика производительности труда" на основе оперативных данных для целей мониторинга хода исполнения Плана мероприятий по обеспечению повышения производительности труда, создания и модернизации высокопроизводительных рабочих мест // Федеральная служба государственной статистики, – 2018. [Электронный ресурс]. – URL: \url{https://www.garant.ru/products/ipo/prime/doc/71786398/}
\bibitem{mineco_prikaz748_28022018_LabourProductivityMethod} Приказ от 28 декабря 2018 г. N~748 Об утверждении методики расчета показателей производительности труда предприятия, отрасли, субъекта Российской Федерации и методики расчета отдельных показателей национального проекта "Производительность труда и поддержка занятости" // Министерство экономического развития Российской Федерации, – 2018. [Электронный ресурс]. – URL: \url{https://economy.gov.ru/material/dokumenty/prikaz_minekonomrazvitiya_rossii_ot_28_dekabrya_2018_g_748.html}
\bibitem{vedomostiRussianIndustryEfficiency2024} Российским промышленникам прописали эффективность / Пресс-служба ФЦК // Ведомости, – 2024. [Электронный ресурс]. – URL: \url{https://www.vedomosti.ru/ideas/development/articles/2024/11/28/1077724-rossiiskim-promishlennikam-propisali-effektivnost}
\bibitem{minecoLabourProductivityProjectPresentation2023} Национальный проект «Производительность труда» (презентационные материалы) // Министерство экономического развития Российской Федерации, – 2023. [Электронный ресурс]. – URL: \url{https://www.economy.gov.ru/material/file/c54c12beb5cac5cf59b2fc0c0f374802/prezentacionnye_materialy.pdf}
\bibitem{tkRf_2001} Трудовой кодекс Российской Федерации // Официальный интернет-портал правовой информации, – 2001. [Электронный ресурс]. – URL: \url{http://pravo.gov.ru/proxy/ips/?docbody=&nd=102074279&rdk=0}
\bibitem{rosstat_postanov1_05012004_FormT13} Постановление от 5 января 2004 г. N~1 Об утверждении унифицированных форм первичной учетной документации по учету труда и его оплаты // Государственный комитет Российской Федерации по статистике, – 2004. [Электронный ресурс]. – URL: \url{https://base.garant.ru/12134807/}
\bibitem{rosstatGdpDeflatorsData2024} Национальные счета, Валовый внутренний продукт, Произведенный ВВП, ВДС годы ОКВЭД2 (с 2011 г.) // Федеральная служба государственной статистики, – 2024. [Электронный ресурс]. – URL: \url{https://rosstat.gov.ru/statistics/accounts}
\bibitem{minecoLabourProductivitySectorProgramsForMay2025} Минэкономразвития: к маю будут разработаны отраслевые программы повышения производительности труда // Министерство экономического развития Российской Федерации, – 2025. [Электронный ресурс]. – URL: \url{https://economy.gov.ru/material/news/minekonomrazvitiya_k_mayu_budut_razrabotany_otraslevye_programmy_povysheniya_proizvoditelnosti_truda.html}
\bibitem{minecoLabourProductivityCoverage2024} Платформа нацпроекта «Производительность труда» по цифровизации бизнеса охватила 3000 предприятий // Министерство экономического развития Российской Федерации, – 2024. [Электронный ресурс]. – URL: \url{https://www.economy.gov.ru/material/news/platforma_nacproekta_proizvoditelnost_truda_po_cifrovizacii_biznesa_ohvatila_3000_predpriyatiy.html}
\bibitem{minecoIntexVolumeExpansion2025} Строительная компания «Интекс» из Магнитогорска на 8\% увеличила выработку производства металлоконструкций // Министерство экономического развития Российской Федерации, – 2025. [Электронный ресурс]. – URL: https://производительность.рф/presscenter/news/stroitelnaya-kompaniya-inteks-iz-magnitogorska-na-8-uvelichila-vyrabotku-proizvodstva-metallokonstrukcij/
\end{thebibliography} 

\end{document}

